\documentclass[11pt,letterpaper,reqno]{amsart}
\usepackage[T1]{fontenc}
\usepackage[utf8]{inputenc}
\usepackage[sc]{mathpazo}
\usepackage[margin=30mm]{geometry}
\usepackage{amsmath,amssymb,amsthm,mathtools,mathrsfs}
\usepackage{microtype}
\usepackage{needspace}
\usepackage[colorlinks=true,linkcolor=blue,citecolor=blue,urlcolor=blue]{hyperref}
\usepackage{bookmark}
\numberwithin{equation}{section}
\newtheorem{theorem}{Theorem}[section]
\newtheorem{proposition}[theorem]{Proposition}
\newtheorem{lemma}[theorem]{Lemma}
\newtheorem{corollary}[theorem]{Corollary}
\theoremstyle{remark}
\newtheorem{remark}[theorem]{Remark}
\DeclareMathOperator{\Var}{Var}
\DeclareMathOperator{\Cov}{Cov}
\DeclareMathOperator{\capD}{cap}
\DeclareMathOperator{\Pois}{Pois}
\DeclareMathOperator{\diag}{diag}
\allowdisplaybreaks[1]
\setlength{\emergencystretch}{2em}
\raggedbottom
\usepackage{etoolbox}
\makeatletter
\patchcmd{\@maketitle}{\global\topskip42\p@\relax}{\global\topskip0\p@\relax}{}{%
  \PackageError{SRW-max-local-time}{Cannot adjust title placement}{Check the amsart version.}}
\makeatother
\hypersetup{pdftitle={Tightness of the maximum local time of planar simple random walk},pdfauthor={Xinyi Li; Yushu Zheng},pdfsubject={Probability theory; random walk; maximum local time}}
\title[Maximum local time of planar simple random walk]{Tightness of the maximum local time of planar simple random walk}
\DeclareRobustCommand{\authormark}[1]{\textsuperscript{\MakeLowercase{#1}}}
\author[X. Li]{Xinyi Li\authormark{1}}
\author[Y. Zheng]{Yushu Zheng\authormark{2}}
\newcommand{\frontaffiliations}{%
  \par\vspace{8pt}
  \begingroup
  \normalfont\footnotesize\centering
  \hypersetup{urlcolor=black}
  \textsuperscript{1}\textit{Beijing International Center for Mathematical Research,
  Peking University}\par
  \href{mailto:xinyili@bicmr.pku.edu.cn}{\textnormal{xinyili@bicmr.pku.edu.cn}}\par\vspace{5pt}
  \textsuperscript{2}\textit{School of Mathematics,
  Shanghai University of Finance and Economics}\par
  \href{mailto:zhengyushu@mail.sufe.edu.cn}{\textnormal{zhengyushu@mail.sufe.edu.cn}}\par
  \endgroup}
\makeatletter
\apptocmd{\@setauthors}{\frontaffiliations}{}{%
  \PackageError{SRW-max-local-time}{Cannot place affiliations on the first page}{Check the amsart version.}}
\makeatother
\date{}
\begin{document}
\begin{abstract}
Let $M_n$ be the largest number of visits to a site by planar simple random walk before step $n$. We prove that $\sqrt{M_n}-(\log n-\log\log n)/\sqrt\pi$ is tight. We also obtain tightness at the first exit from a disk and quantitative, non-optimal tail bounds. The proof compares complete excursion records with a critical geometric Galton--Watson process. For the lower bound, we count sites with high occupation whose nested equilibrium averages satisfy a continuous barrier. A Gaussian density comparison under exponential tilting gives the one-point estimate, while the common average path controls nearby pairs. A one-sided Gaussian comparison gives the upper bound. Successive disk exits amplify the lower bound, and an exit-time comparison yields the result at deterministic times.
\end{abstract}
\maketitle

\section{Introduction}\label{sec:intro}
The maximum local time of a random walk measures the greatest concentration of visits along its path. For recurrent walks, repeated returns produce large local times, while nearby sites share much of the same excursion history. The leading order of the maximum and its fluctuations therefore involve different aspects of this dependence.

We consider nearest-neighbor simple random walk on the planar lattice. Its local time at a site is the number of visits to that site, and we study the largest such count before a deterministic time or before exit from a disk. Dembo, Peres, Rosen and Zeitouni~\cite{DPRZ} proved the Erd\H{o}s--Taylor conjecture: the maximum before time $n$, divided by $(\log n)^2$, converges almost surely to $1/\pi$. Rosen~\cite[Theorem 1.1]{ET} subsequently gave a proof based directly on random-walk excursions.

The next question is to determine the logarithmic correction to the square root of the maximum and to control the remaining fluctuations. For planar Brownian motion, Rosen~\cite[Theorem 1.1]{GW} proved tightness for the corresponding maximum of normalized small-disk occupation times and raised the discrete question in~\cite[Section 1.2]{GW}. The Brownian observable averages occupation over a disk. The lattice problem requires control of visits to an individual site at the same fluctuation scale.

In this paper we prove tightness after the logarithmic correction, both at disk exit and at deterministic times, and obtain quantitative bounds for the two tails. The proof retains the discrete excursion records down to their terminal occupations. For the lower bound, constraints on nested equilibrium averages allow us to count individual high sites and control their dependence. A comparison with the discrete Gaussian free field supplies the upper bound. The main theorems concern the unrestricted maximum of the original walk.

\subsection{Model and main results}\label{sec:results}
Let $S=(S_k)_{k\ge0}$ be simple random walk on $\mathbb Z^2$, with $S_0=0$ and independent increments uniformly distributed over the four nearest-neighbor vectors. We write $\mathbb P_x$ and $\mathbb E_x$ for its law and expectation when started at $x$. All logarithms are natural, and $|\cdot|$ denotes the Euclidean norm.

For each integer $n\ge3$, set
\begin{equation}\label{eq:1.1}
\ell_n(x)=\sum_{k=0}^{n-1}\mathbf1_{\{S_k=x\}},\qquad
M_n^{\mathrm{time}}=\max_{x\in\mathbb Z^2}\ell_n(x),\qquad
\beta_n=\frac{\log n-\log\log n}{\sqrt\pi}.
\end{equation}
The half-open time interval in this definition will also be used for stopped local times.

\begin{theorem}\label{thm:time}
The family $\sqrt{M_n^{\mathrm{time}}}-\beta_n$, indexed by all integers $n\ge3$, is tight. Equivalently,
\begin{equation}\label{eq:1.2}
\lim_{K\to\infty}\limsup_{n\to\infty}
\mathbb P_0\!\left(\left|\sqrt{M_n^{\mathrm{time}}}-\beta_n\right|>K\right)=0.
\end{equation}
The same conclusion, with the same centering, holds for
$\widehat M_n=\max_x\sum_{k=0}^{n}\mathbf1_{\{S_k=x\}}$.
\end{theorem}

On the same infinite path, for a real number $R>1$, define
\begin{equation}\label{eq:1.3}
\begin{gathered}
D_R=\{x\in\mathbb Z^2:|x|<R\},\qquad
\tau_R=\inf\{k\ge0:S_k\notin D_R\},\\
L_R(x)=\sum_{k<\tau_R}\mathbf1_{\{S_k=x\}},\qquad
M_R=\max_{x\in\mathbb Z^2}L_R(x),\qquad
b_R=\frac{2\log R-\log\log R}{\sqrt\pi}.
\end{gathered}
\end{equation}
We first prove a spatial version of Theorem~\ref{thm:time}.

\begin{theorem}\label{thm:disk}
As $R$ ranges over all sufficiently large real numbers,
\begin{equation}\label{eq:1.4}
\lim_{K\to\infty}\limsup_{R\to\infty}
\mathbb P_0\!\left(\left|\sqrt{M_R}-b_R\right|>K\right)=0.
\end{equation}
Moreover, there are constants $C,c_*,s_*>0$ such that, for every fixed $s\ge s_*$,
\begin{equation}\label{eq:1.6}
\limsup_{R\to\infty}\mathbb P_0(\sqrt{M_R}>b_R+s)
\le C(1+s)e^{-3s/2},
\end{equation}
and
\begin{equation}\label{eq:1.7}
\limsup_{R\to\infty}\mathbb P_0(\sqrt{M_R}<b_R-s)
\le\exp\{-c_*e^{\sqrt\pi s/4}\}.
\end{equation}
\end{theorem}

The maximum in Theorem~\ref{thm:disk} is taken over the entire disk. We do not claim optimal tail exponents. For any fixed $R_0>1$, the theorem also gives tightness of the family indexed by $R\ge R_0$: monotonicity on the same path handles each bounded interval of radii.

\subsection{Proof outline}\label{sec:outline}
The proof begins with two estimates for the disk maximum: a decaying upper tail and a probability, bounded away from zero, of reaching the centering level $b_R$. Successive exits from smaller disks provide independent chances to reach this level with a bounded shift in the centering. This amplifies the positive probability into a lower-tail bound. Comparing the walk at time $n$ with its exits from disks of radii proportional to $\sqrt n$ then gives the deterministic-time theorem.

To prove the spatial estimates, we record the ordered inward excursions through concentric annuli, together with the occupations generated by their terminal segments. The full record is needed because constraints on excursion counts must be imposed jointly with high occupation at the observed sites. Equilibrium entrance laws give reference marks whose mean occupation is constant across those sites. A multiplicative comparison of records reduces the radial counts to a critical geometric Galton--Watson process without discarding the terminal occupations.

For the upper bound, a one-sided Gaussian comparison controls the maximum in each terminal block. The discrete Gaussian free field maximum estimate, combined with a bound on its harmonic component, gives the required terminal probability. We then sum over the coarse terminal counts and the random number of returns to the root annulus. A separate estimate for spatial barrier violations ensures that this argument covers all high sites.

For the lower bound, the observation vector consists of a candidate site and its nested equilibrium averages. The covariance is Brownian covariance minus a summable diagonal correction. A multiplicative Gaussian density comparison preserves the probability of the rare barrier event even when the smallest averaging radius is fixed. Independent auxiliary Brownian bridges make the average constraint continuous; after fixing these bridges, the constraint remains convex in the observed occupations. Exponential tilting and convex-set normal approximation give a terminal probability whose prefactor is balanced by the sum over coarse terminal counts.

The second moment depends on the scale at which two sites separate. For nearby sites, retaining the entire common average path gives a killed Brownian kernel and a summable distance bound. For earlier separations, we condition on the exterior excursion record. An exact identity for the retained trace removes the prescribed exit conditions, after which a second comparison returns the calculation to the original joint law of the two root counts. These estimates give the positive probability needed for the lower-tail amplification.

\subsection{Further related work}\label{sec:related}
Excursion barriers connect planar thick points with branching processes. In Rosen's Brownian proof~\cite{GW}, barrier estimates for a critical geometric Galton--Watson process control radial excursion counts and separate the first- and second-moment calculations. We use the version stated in~\cite[Theorem 9.1]{GW}, with its parameter conditions checked at the points of application. The comparison with these counts is carried out on the discrete walk's retained records, including their terminal marks.

The discrete Gaussian free field provides another description of logarithmic spatial correlations. Bramson, Ding and Zeitouni~\cite{BDZ} proved convergence in law of its centered maximum. Our upper bound uses their maximum tail estimate and a harmonic-field increment estimate, with the covariance normalized by the SRW Green function. The occupation marks themselves are not Gaussian: the upper bound uses a one-sided comparison, whereas the lower bound uses normal approximation under a fixed-number mark tilt. The remaining potential-theoretic and normal-approximation inputs are collected in Sections~\ref{sec:records} and~\ref{sec:tools}.

\subsection{Organization}\label{sec:organization}
Section~\ref{sec:mainproofs} derives the main theorems from the two spatial estimates. Section~\ref{sec:records} constructs and compares the excursion records, and Section~\ref{sec:tools} collects the branching-process and Gaussian inputs. Section~\ref{sec:upper} proves the upper bound. Section~\ref{sec:events} defines the lower-bound event and proves its first-moment estimate. Sections~\ref{sec:tworoot} and~\ref{sec:second} establish the pair estimates and the lower bound by the second-moment method.

\subsection{AI-assisted proof development and verification status}
\label{subsec:ai-declaration}

Generative AI tools were used to explore proof strategies, generate the mathematical proof draft, prepare the manuscript, and develop the accompanying Lean~4 formalization. The authors take responsibility for the mathematical content, its verification, and its exposition.

The accompanying Lean~4 formalization proves the registered main tightness and tail claims conditional on six explicit literature input families; see Appendix~\ref{sec:formalization}. This formal acceptance does not constitute an independent human verification of every argument in the manuscript. The checksum-bound source is stored in \texttt{lean/source/}\allowbreak\texttt{SRW\_max\_local\_time.tex}; the publication version retains its mathematical body, with different front matter and this status statement. Further polishing of the manuscript and its exposition is planned.

\section{Proofs of Theorems~\ref{thm:disk} and~\ref{thm:time}}\label{sec:mainproofs}
The following two propositions contain the spatial estimates needed for both theorems. The upper bound is proved in Section~\ref{sec:upper}. The positive-probability lower bound is proved in Sections~\ref{sec:events}--\ref{sec:second}. We first derive the main theorems from them.

\begin{proposition}\label{prop:upper}
There are constants $C,s_*>0$ such that, for each fixed $s\ge s_*$,
\[
\limsup_{R\to\infty}\mathbb P_0(\sqrt{M_R}>b_R+s)
\le C(1+s)e^{-3s/2}.
\]
\end{proposition}

\begin{proposition}\label{prop:seed}
There are $p_*>0$ and $R_*<\infty$ such that
\begin{equation}\label{eq:1.5}
\mathbb P_0(\sqrt{M_R}\ge b_R)\ge p_*,\qquad R\ge R_*.
\end{equation}
\end{proposition}

\subsection{An exit-time estimate}
We will use the following estimate both to justify the successive disk exits and to compare exit times with deterministic times.
\begin{lemma}\label{lem:exit}
For $a>1$, the exit time $\tau_a$ is integrable and
\begin{equation}\label{eq:2.1}
a^2\le\mathbb E_0\tau_a=\mathbb E_0|S_{\tau_a}|^2\le(a+1)^2.
\end{equation}
For $n\ge1$ one also has
\begin{equation}\label{eq:2.2}
\mathbb P_0(\tau_a\le n)\le n/a^2.
\end{equation}
\end{lemma}
\begin{proof}
Bounded optional stopping for $|S_k|^2-k$ at $T_m=\tau_a\wedge m$ gives
\begin{equation}\label{eq:2.3}
\mathbb E_0T_m=\mathbb E_0|S_{T_m}|^2\le(a+1)^2.
\end{equation}
Monotone convergence proves integrability; bounded overshoot then gives~\eqref{eq:2.1} by dominated convergence. Applying the same identity at $\tau_a\wedge n$ gives $a^2\mathbb P_0(\tau_a\le n)\le\mathbb E_0(\tau_a\wedge n)\le n$.
\end{proof}
Translation gives the same statements for disks centered at the starting point. In particular, exit times from all finite domains are integrable.

\subsection{Proof of Theorem~\ref{thm:disk}}
Successive exits from disks centered at the current position give independent path segments. Applying Proposition~\ref{prop:seed} to each segment will yield the lower tail.
\begin{proof}
Write $\mathfrak c=2/\sqrt\pi$. Fix $a>1$ and an integer $J$ with $J(a+1)<R$, and define
\begin{equation}\label{eq:10.1}
\sigma_0=0,\qquad
\sigma_{i+1}=\inf\{k\ge\sigma_i:|S_k-S_{\sigma_i}|\ge a\}.
\end{equation}
The first $J$ segments lie in $D_R$, since their displacement from the origin is at most $J(a+1)<R$. Their translated paths on the half-open intervals $[\sigma_i,\sigma_{i+1})$ are independent copies of a walk killed on leaving $D_a$. Each segment's visit counts are bounded by those before $\tau_R$, whether or not the spatial ranges of the segments overlap. Consequently, for every integer $k\ge1$,
\begin{equation}\label{eq:10.2}
\mathbb P_0(M_R<k)\le\mathbb P_0(M_a<k)^J.
\end{equation}

For sufficiently large fixed $s$, take
\begin{equation}\label{eq:10.3}
Q=\lfloor e^{s/(2\mathfrak c)}\rfloor,\qquad
a=R/Q,\qquad J=\lfloor Q/2\rfloor.
\end{equation}
As $R\to\infty$ these choices eventually satisfy the required geometric inequality, and
\begin{equation}\label{eq:10.4}
b_R-b_a=\mathfrak c\log Q+
\pi^{-1/2}\log\left(1-\frac{\log Q}{\log R}\right)\le s/2.
\end{equation}
Thus $b_a\ge b_R-s$. Apply~\eqref{eq:10.2} with $k=\lceil b_a^2\rceil$ and use Proposition~\ref{prop:seed} to obtain
\begin{equation}\label{eq:10.5}
\limsup_{R\to\infty}\mathbb P_0(\sqrt{M_R}<b_R-s)
\le(1-p_*)^{\lfloor Q/2\rfloor}
\le\exp\{-c_*e^{\sqrt\pi s/4}\}.
\end{equation}
Proposition~\ref{prop:upper} proves the other tail. Their sum tends to zero as $s\to\infty$, proving tightness. The estimates hold for all sufficiently large real radii. On a fixed compact subinterval of $(1,\infty)$, monotonicity bounds the maximum by the maximum at the largest radius, while the center is bounded.
\end{proof}

\subsection{Proof of Theorem~\ref{thm:time}}
We place time $n$ between the exits from two disks on the same path. Their radii differ from $\sqrt n$ by fixed factors, so their centerings differ from $\beta_n$ by bounded amounts.
\begin{proof}
Fix $A>1$, set $\rho_-=\sqrt n/A$ and $\rho_+=A\sqrt n$, and let
$E_{n,A}=\{\tau_{\rho_-}\le n\le\tau_{\rho_+}\}$. Lemma~\ref{lem:exit} and Markov's inequality give
\begin{equation}\label{eq:11.2}
\mathbb P_0(E_{n,A}^c)\le2A^{-2}+\frac{2}{A\sqrt n}+\frac1n.
\end{equation}
On $E_{n,A}$ the visit intervals are nested, so
\begin{equation}\label{eq:11.3}
M_{\rho_-}\le M_n^{\mathrm{time}}\le M_{\rho_+}.
\end{equation}
Let $\kappa_0=(\log2)/\sqrt\pi$ and $\mathfrak c=2/\sqrt\pi$. Direct substitution yields, for fixed $A$,
\begin{equation}\label{eq:11.4}
b_{\rho_\pm}-\beta_n
=\kappa_0\pm\mathfrak c\log A
-\pi^{-1/2}\log\left(1\pm\frac{2\log A}{\log n}\right)
=\kappa_0\pm\mathfrak c\log A+o(1).
\end{equation}
Put $X_n=\sqrt{M_n^{\mathrm{time}}}-\beta_n$ and $C_A=\kappa_0+\mathfrak c\log A+1$. For all sufficiently large $n$,~\eqref{eq:11.3}--\eqref{eq:11.4} imply
\[
\mathbb P_0(|X_n|>s)
\le\mathbb P_0(E_{n,A}^c)
+\sum_{\epsilon\in\{-,+\}}
\mathbb P_0\bigl(|\sqrt{M_{\rho_\epsilon}}-b_{\rho_\epsilon}|>s-C_A\bigr).
\]
For each fixed $A$, disk tightness makes the last two limsup probabilities tend to zero as $s\to\infty$. Letting $A\to\infty$ in~\eqref{eq:11.2} proves tightness at deterministic times. The finite initial range of $n$ is bounded by $M_n^{\mathrm{time}}\le n$. Finally $\widehat M_n=M_{n+1}^{\mathrm{time}}$ and $\beta_{n+1}-\beta_n\to0$, proving the stated convention with the last visit included.
\end{proof}

The same comparisons also give the quantitative estimate
\begin{equation}\label{eq:11.9}
\begin{split}
\limsup_{n\to\infty}\mathbb P_0(|X_n|>s)
\le{}&2e^{-s/(2\mathfrak c)}+C(1+s)e^{-3s/4}\\
&+\exp\{-c_*e^{\sqrt\pi s/8}\}
\end{split}
\end{equation}
for sufficiently large fixed $s$. Indeed, take $A=e^{s/(4\mathfrak c)}$ before sending $n\to\infty$. The upper and lower disk deviations required by~\eqref{eq:11.3} have sizes at least $s/2$ for large $s$, by~\eqref{eq:11.4}. The disk tails and~\eqref{eq:11.2} then give~\eqref{eq:11.9}.

\section{Excursion records and equilibrium marks}\label{sec:records}
Both spatial estimates use a decomposition of local time into visits made during nested excursions. We first collect the potential-theoretic estimates needed for this decomposition, then construct an ordered record whose law can be compared multiplicatively with a geometric Galton--Watson process. Keeping every visit in the record will allow us to apply this comparison directly to local-time events. Throughout the spatial argument,
\[
g=\frac2\pi,
\]
and $c,C$ denote positive constants that may change from line to line. They are independent of $R$; dependence on fixed geometric parameters will be indicated when relevant. For a lattice center $b$, write $B(b,a)=\{z\in\mathbb Z^2:|z-b|<a\}$.

\subsection{Green functions and the one-point tail}\label{sec:green}
For a finite lattice domain $D$, let
\[
P_D(x,y)=\tfrac14\mathbf1_{\{|x-y|=1\}},\qquad x,y\in D,
\quad G_D=(I-P_D)^{-1}.
\]
The inverse exists by Lemma~\ref{lem:exit}. Write $\tau_D$ for the first exit time from $D$, and let $H_K$ and $H_K^+$ be the first hitting times of $K$ at times at least zero and at least one, respectively. When $K\subset D$, $G_K$ denotes the restriction of $G_D$ to $K\times K$, not a Green function killed on leaving $K$.

Define the equilibrium measure and capacity relative to $D$ by
\[
e_K^D(x)=\mathbf1_K(x)\mathbb P_x(\tau_D<H_K^+),\qquad
\operatorname{cap}_D(K)=\sum_x e_K^D(x).
\]
A last-visit decomposition gives
\begin{equation}\label{eq:2.4}
\sum_{x\in K}G_D(z,x)e_K^D(x)=\mathbb P_z(H_K<\tau_D).
\end{equation}
By symmetry, $G_Ke_K=\mathbf1$ and $e_K^TG_K=\mathbf1^T$ on $K$. Superscripts specifying the killing domain will be omitted when that domain is fixed.

For the present walk, the potential kernel satisfies
\begin{equation}\label{eq:2.5}
a(z)=g\log|z|+\kappa+O(|z|^{-2}),\qquad
\kappa=(2\gamma_E+3\log2)/\pi,
\end{equation}
where $\gamma_E$ is Euler's constant; see~\cite[Theorem 4.4.4]{LL}. Since $(P-I)a=\mathbf1_{\{0\}}$, stopping in a bounded domain gives
\begin{equation}\label{eq:2.6}
G_D(x,y)=\mathbb E_xa(S_{\tau_D}-y)-a(x-y).
\end{equation}
In particular, the one-step overshoot implies
\begin{equation}\label{eq:2.7}
G_{B(b,r)}(b,b)=g\log r+\kappa+O(r^{-1}).
\end{equation}
For $x$ in any fixed proportional interior of $D_R$, one similarly has $G_{D_R}(x,x)=g\log R+O(1)$, while
\begin{equation}\label{eq:2.8}
G_{D_R}(0,x)\le C\left(1+\log\frac R{|x|\vee1}\right).
\end{equation}

After its first visit to $x$, the number of visits is geometric, with probability $G_D(x,x)^{-1}$ of no further return. Therefore, for integers $k\ge1$,
\begin{equation}\label{eq:2.9}
\mathbb P_0(L_R(x)\ge k)=
\frac{G_{D_R}(0,x)}{G_{D_R}(x,x)}
\left(1-\frac1{G_{D_R}(x,x)}\right)^{k-1}.
\end{equation}
We will use the hitting factor in this identity when summing the upper tail over sites near the origin.

\subsection{Annulus probabilities and equilibrium averages}\label{sec:annuli}
Stopping the potential kernel in the annulus $A<|x-b|<B$, and using the bounded overshoot, gives
\begin{equation}\label{eq:2.10}
\mathbb P_x(\tau_{B(b,B)}<H_{B(b,A)})
=\frac{\log(|x-b|/A)+O(A^{-1})}
       {\log(B/A)+O(A^{-1})}.
\end{equation}
For a relative estimate, the error must be divided by the logarithmic distance in the numerator. In each application this distance is either bounded below or retained in the error bound.

Let $K=B(b,a)$, $a\ll m$, and $D=B(b,em+H)$, where $H=o(m)$. Put
$\omega=e_K^D/\operatorname{cap}_D(K)$ and $V=1/\operatorname{cap}_D(K)$.
Evaluating~\eqref{eq:2.4} at the center, where $G_D(b,x)=g\log(m/a)+O(1)$ uniformly on the support of $e_K^D$, yields
\begin{equation}\label{eq:2.11}
V/g=\log m-\log a+O(1).
\end{equation}
At a boundary point of $K$, first take the step out of $K$ and then use~\eqref{eq:2.10}. This gives
\[
e_K^D(x)\le\frac{C}{a\log(m/a)}.
\]
After division by the capacity,
\begin{equation}\label{eq:2.12}
\max_x\omega(x)\le C/a,\qquad \|\omega\|_2^2\le C/a.
\end{equation}
For a square of side comparable to $a$, with the same separation from the killing boundary, the energy estimate~\eqref{eq:2.11} remains valid. For such squares we will only use the elementary bound $\|\omega\|_2^2\le1$.

We also use the positive harmonic difference estimate and Harnack inequality in~\cite[Theorems 6.3.8(b) and 6.3.9]{LL}. If $f\ge0$ is harmonic in $B(b,B)$, then, for $|x-b|\le2A\ll B$,
\begin{equation}\label{eq:2.13}
f(x)=f(b)\{1+O(A/B)\}.
\end{equation}
Because the error is relative to $f(b)$, this estimate also controls the harmonic measure of a specified exit point, uniformly over the choice of exit point.

\subsection{Construction of reference marks and entrance-law comparison}\label{sec:marks}
We next choose a reference law for the visits made after an entrance into the terminal annulus. Its equilibrium starting measure makes the mean occupation constant over the observation set. Let $K\subset B(b,a)$, $0<a\le m/2$, $0\le H\le m/8$, and $D=B(b,em+H)$, with $m$ sufficiently large. Set
\[
\alpha=g\log\frac{em+H}{m-H}.
\]
If $\alpha\operatorname{cap}_D(K)\le1$, define a reference mark by taking a killed random-walk path with subprobability law
\begin{equation}\label{eq:2.14}
\alpha\sum_{v\in K}e_K^D(v)\mathbb P_v^D,
\end{equation}
and adding a zero mark of mass $1-\alpha\operatorname{cap}_D(K)$. We retain only visits to $K$, or a specified projected record after the first entrance into $K$. By~\eqref{eq:2.4}, the expected number of visits to each $z\in K$ is exactly $\alpha$. Enlarging $K$ and then restricting the observation gives the same mark law on the smaller set, by the first-hitting decomposition and symmetry.

For a true entrance point $z$ with $m-H-1\le|z-b|\le m-H+1$, let $h_z(v)$ denote the probability of first entering $K$ at $v$ before exit from $D$. We first establish a uniform comparison that does not require $a/m$ to tend to zero.

Put $L=B(b,m-H)$. For each $v\in K$, the function $x\mapsto h_x(v)$ is nonnegative and harmonic in $D\setminus K$, with value zero outside $D$. Any two points of the entrance layer can be joined by a bounded number of overlapping balls of radius proportional to $m$, whose fixed enlargements lie in $D\setminus B(b,m/2)$. The Harnack inequality from Section~\ref{sec:annuli} along this chain gives $h_x(v)\asymp h_z(v)$ throughout the layer, uniformly also when the entrance probability is small. The equilibrium measure of $L$ is supported on its inner unit boundary layer, so this comparison applies on its support.

The equilibrium sweeping identity is
\[
\sum_{x\in L}e_L^D(x)h_x(v)=e_K^D(v).
\]
Indeed, write $h_\cdot(v)=G_D\nu$, where $\nu=(I-P_D)h_\cdot(v)$ is supported on $K$. Summing against either $e_L^D$ or $e_K^D$ gives $\sum_{y\in K}\nu(y)$ by~\eqref{eq:2.4} and symmetry. The latter sum against $e_K^D$ equals $e_K^D(v)$, since $h_y(v)=\mathbf1_{\{y=v\}}$ for $y\in K$.
Also,~\eqref{eq:2.6} gives $G_D(b,x)=\alpha+O(m^{-1})$ on the support of $e_L^D$. Evaluating~\eqref{eq:2.4} at $b$ therefore yields
\[
\operatorname{cap}_D(L)^{-1}=\alpha+O(m^{-1}),\qquad
\tfrac12\le\alpha\operatorname{cap}_D(L)\le2.
\]
Combining sweeping with the Harnack comparison gives fixed constants $c_0,C_0>0$ such that
\[
c_0\alpha e_K^D(v)\le h_z(v)\le C_0\alpha e_K^D(v),
\qquad v\in K.
\]
All constants here are independent of $K,a,m,H,z$ in the stated range. The multiplicative inequalities also cover zero equilibrium charges.

For the refined relative estimate, suppose additionally that $2a<m-H-2$. Uniformly over all positive entrance atoms,
\begin{equation}\label{eq:2.15}
h_z(v)=\alpha e_K^D(v)(1+O(\epsilon)),\quad
\epsilon=C\frac{a+1}{m}\log\frac ma,
\end{equation}
First suppose $a\le m/4$ and put $A=a\vee a_0$, where $a_0\ge1$ is a fixed sufficiently large radius for~\eqref{eq:2.10}. We take $m$ large enough that $a_0\le m/4$. Let $f(v)=G_D(z,v)$. By~\eqref{eq:2.6}, $f(v)=\alpha+O((a+1)/m)$ on the outer unit boundary layer of $B(b,2A)$; the fixed $a_0$ is absorbed into the constant. Let $q(v)=\mathbb P_v(\tau_D<H_K)$. There, avoiding the entire $A$-ball gives, by~\eqref{eq:2.10}, the lower bound $q(v)\ge c/\log(m/A)\ge c/\log(m/a)$. First hitting and symmetry imply, on that outer boundary,
\[
G_{D\setminus K}(z,v)
=f(v)-\mathbb E_v[f(S_{H_K});H_K<\tau_D]
=\alpha q(v)+O((a+1)/m).
\]
Both functions vanish on $K$ and are harmonic in the punctured ball. Applying the maximum principle to their difference propagates the relative estimate to every exterior neighbor of $K$. Summing the last step into $K$ proves~\eqref{eq:2.15} in this range. If $m/4<a\le m/2$, the uniform atom comparison above gives the same relative estimate after increasing the constant, because
\[
\frac{a+1}{m}\log\frac ma\ge\frac{\log2}{4}.
\]
In particular this proves~\eqref{eq:2.15} throughout its stated separated range without requiring its error to be small.

After the first entrance into $K$, the true and reference paths have the same killed transition law. Thus~\eqref{eq:2.15} compares the entire retained record with an error independent of its number of visits. If the reference zero mass is at least $1/2$, the zero atom has the same relative comparison, as detailed in Section~\ref{sec:recordcomparison}. The uniform comparison also extends to the zero atom: by~\eqref{eq:2.10}, avoiding $B(b,m/2)$ until exit from $D$ has probability bounded below by a fixed positive constant from the entrance layer. Hence the actual zero mass is uniformly positive. If $\alpha\operatorname{cap}_D(K)\le1/2$, the reference zero mass is at least $1/2$, which supplies the upper comparison. Consequently the full actual and reference mark laws, denoted by $P_z$ and $Q$, satisfy
\[
c\,Q(E)\le P_z(E)\le C\,Q(E)
\]
for every mark event $E$, uniformly for $0<a\le m/2$. These inequalities survive every specified projection of the retained record. The pre-entrance path is not part of this observation; conditioning on a specified final exit point will be treated separately in Section~\ref{sec:tworoot}.

\subsection{Generating functions and moments of discrete marks}\label{sec:markformulas}
The generating function of a reference mark provides both the Gaussian comparison for the upper bound and the covariance needed for the lower bound. For a finite observation set $A\subset K$, let $G=G_D|_{A\times A}$, $Q=\operatorname{diag}(v_z-1)$, and let $Y$ be the visit vector of one reference mark. Then
\begin{equation}\label{eq:2.16}
\mathbb E\prod_{z\in A}v_z^{Y_z}
=1+\alpha\mathbf1^TQ[I-(G-I)Q]^{-1}\mathbf1.
\end{equation}
To verify the formula near $v_z=1$, expand the product over selected visits and order their times strictly increasingly. The intensity of the first selected visit is $\alpha$, and the kernel between two strictly different visit times is $G-I$. Summing the matrix geometric series gives~\eqref{eq:2.16}. The identity extends to $0\le v_z\le1$ by the finite-chain resolvent formula. Positive parameters will only be used within an explicitly controlled convergence range.

Differentiation yields
\begin{equation}\label{eq:2.17}
\mathbb EY=\alpha\mathbf1,\qquad
\mathbb EYY^T=2\alpha G-\alpha I,\qquad
\operatorname{Cov}Y=2\alpha G-\alpha I-\alpha^2\mathbf1\mathbf1^T.
\end{equation}
Each coordinate has the zero--geometric mixture law
\begin{equation}\label{eq:2.18}
\mathbb P(Y_z=0)=1-\alpha/G_{zz},\qquad
\mathbb P(Y_z=l)=\alpha G_{zz}^{-2}(1-G_{zz}^{-1})^{l-1},\quad l\ge1.
\end{equation}
Consequently, if $G_{zz}\le Cd$ and $\alpha\asymp1$, then $\mathbb E Y_z^p\le C_pd^{p-1}$. Nonnegative averages with weights summing to one satisfy the same bound by Jensen's inequality. For signed combinations whose coefficient absolute sum is bounded, geometric exponential moments and H\"older's inequality give a third centered moment bounded by $Cd^3$ under single-mark tilts of order $t^{-1}$ as $t\to\infty$, provided $d/t\to0$. This estimate also bounds the Taylor remainders used below.

For $k=u^2/2$ independent marks, let $H=\sum_{i=1}^kY_i$ and put
\[
F=\frac{H-\alpha k\mathbf1}{u\sqrt{\alpha g}}.
\]
The exact covariance is
\begin{equation}\label{eq:2.19}
\operatorname{Cov}F=G/g-I/(2g)-\alpha\mathbf1\mathbf1^T/(2g).
\end{equation}
The two correction terms in~\eqref{eq:2.19} account for discrete visit times and the fixed number of marks; both enter the lower-bound calculation.

\subsection{Decomposition into complete excursions}\label{sec:physical}
We now define the nested excursions that assign each visit to a terminal mark. Given scales $m,r$, set
\begin{equation}\label{eq:2.20}
\ell=\left\lfloor\log\frac R{32m}\right\rfloor,\quad
\rho_j=me^{\ell-j},\quad N_R=\ell-1,\quad
H_0=4r+10,\quad A_j=\rho_j-H_0,\quad B_j=\rho_j+H_0.
\end{equation}
We only use parameters with $m\gg H_0$ and $N_R\to\infty$. Fix the center $b$.
A root starts at the first entrance into $B(b,A_1)$ and ends at the subsequent exit from $B(b,B_0)$. After each such exit, the next return to $B(b,A_1)$ before $\tau_R$ starts another root. Let $J=N_0$ be their total number. For a root starting at level $k$, replace the indices $0,1$ by $k,k+1$ and write $N_k=J$.

A parent at level $j-1$ runs from entrance at radius $A_j$ to exit at radius $B_{j-1}$. Inside it, enumerate in chronological order all segments from entrance at radius $A_{j+1}$ to exit at radius $B_j$; these are its children. Their total over all parents is $N_j$. Nesting and nearest-neighbor steps ensure that every started child is completed before its parent ends. The terminal, level-$N_R$ segments run from radius $m-H_0$ to radius $em+H_0$.

Suppose the original starting point is outside the root entrance ball, all root outer balls lie in $D_R$, and $K\subset B(b,m-H_0-2)$. If $Y_i$ now denotes the actual visit vector of terminal segment $i$, then pathwise
\begin{equation}\label{eq:2.21}
(L_R(z):z\in K)=\sum_{i=1}^{N_{N_R}}Y_i.
\end{equation}
Indeed, there are no visits to $K$ before the first root, between roots, or after the last root. Within a parent, the portions between children contain no visits to the more interior observation set. Iterating this observation assigns every visit to a unique terminal segment. All segments are half-open, so a common endpoint is counted only as the beginning of the later segment. Section~\ref{sec:recordcomparison} compares the joint law of these actual marks with independent reference marks.

For the upper bound, we also need traversal counts centered at individual observation points. These counts will dominate the excursion counts of any nearby block. For an observation point $x$, define the inward records at level $j$ by successive entrance and exit times. Suppressing $j,x$ in the stopping-time notation, put
\[
\sigma_1=\inf\{n\ge0:|S_n-x|<\rho_{j+1}\},
\]
and, recursively,
\begin{equation}\label{eq:traversal-stopping}
\begin{aligned}
\eta_i&=\inf\{n\ge\sigma_i:|S_n-x|\ge\rho_j\},\\
\sigma_{i+1}&=\inf\{n\ge\eta_i:|S_n-x|<\rho_{j+1}\}.
\end{aligned}
\end{equation}
An infimum over the empty set is infinity, and
$\mathcal T_j^x=\sum_{i\ge1}\mathbf1_{\{\sigma_i<\tau_R\}}$.
Apart from an entrance at time zero, nearest-neighbor steps place the entrance and exit points in unit-thickness boundary layers. This upper-bound convention includes the initial inward piece if the walk starts between the two shells; if it starts in the inner ball, $\sigma_1=0$. Every later record requires an intervening outer-shell visit.

Suppose $|x-b|\le r$. At the entrance of a block child at radius $A_{j+1}$, the path lies in $B(x,\rho_{j+1})$; at its exit at radius $B_j$, it lies outside $B(x,\rho_j)$, since $H_0>r+3$. Assign the child to the unique interval $[\sigma_i,\eta_i)$ containing its entrance time. Consecutive children are assigned distinct intervals because the first child's exit occurs after $\eta_i$. The same separation holds across different parents, whose exits are farther out. Consequently,
\begin{equation}\label{eq:2.22}
N_j\le\mathcal T_j^x.
\end{equation}
Thus a barrier for the point-centered traversal counts also bounds the corresponding block counts.

\subsection{Multiplicative comparison under count bounds}\label{sec:recordcomparison}
The comparison keeps the ordered forest of complete segments together with the projected mark attached to each terminal segment. For fixed radii and depth, these records form a countable space: the forest is finite, and each mark is the image of a finite lattice path word. We equip this space with the sigma-field generated by its individual records. Entrance and exit positions used in the construction are integrated out unless they belong to the specified terminal projection.

In a parent segment, the next entrance into the inner ball competes with exit from the outer ball. At the initial middle entrance and after each middle return,~\eqref{eq:2.10} gives probabilities
\begin{equation}\label{eq:2.23}
\frac12+O(H_0/\rho_{j+1}),
\end{equation}
uniformly in the actual position. The reference offspring law is therefore the critical geometric law $\mathbb P(\xi=l)=2^{-(l+1)}$, $l\ge0$.

At the terminal level, first consider the complete mark before any further projection. Assume explicitly that its reference positive mass is $q=\alpha\operatorname{cap}_D(K)\le1/2$, as will hold in the applications below. In the separated range,~\eqref{eq:2.15} compares every positive complete mark. To include the zero mark, let $\delta$ bound the relative entrance-law error. If $Q$ is the reference mark law and $P_z$ the actual law from an admissible entrance $z$, then
\[
|P_z(0)-Q(0)|\le\delta q,\qquad
\left|\frac{P_z(0)}{Q(0)}-1\right|
\le\frac{\delta q}{1-q}.
\]
Hence the same relative error applies to the full terminal law, uniformly over admissible entrances. To pass to a two-sided exponential comparison, write $\epsilon_0=(a+1)m^{-1}\log(m/a)$. For $a\le m/4$, the separation $2a<m-H-2$ holds for large $m$, and $|P_z(E)-Q(E)|\le\delta Q(E)$ with $\delta=C_1\epsilon_0$. If $\delta\le1/2$, use $1-\delta\ge e^{-2\delta}$ and $1+\delta\le e^\delta$. If $\delta>1/2$, the fixed positive lower comparison from Section~\ref{sec:marks} gives the lower exponential bound after increasing the constant, since $\epsilon_0>1/(2C_1)$; the upper bound still follows from $1+\delta\le e^\delta$. For $m/4<a\le m/2$, use the fixed two-sided comparison directly and the bound $\epsilon_0\ge(\log2)/4$. We have therefore proved, throughout the full half-radius range,
\[
e^{-C\epsilon_0}Q(E)\le P_z(E)\le e^{C\epsilon_0}Q(E).
\]
This comparison includes the zero mark and does not assume that the terminal relative error is small.

We propagate the terminal comparison backwards through the tree, integrating each exit position before applying the comparison for the preceding child. At level $j$ write $U_j(x,dz)$ for the kernel of entering $B(b,A_{j+1})$ before exiting $B(b,B_{j-1})$, and $V_j(x,dy)$ for the competing exit kernel. Their total masses satisfy~\eqref{eq:2.23} both at the initial position on the $A_j$ boundary and at a return position on the $B_j$ boundary. Let $S_j(z;r,d\xi)$ be the joint kernel of a child record $r$ and its exit position $\xi$. Suppose inductively that its marginal is bounded above and below by $e^{\pm\eta(r)}Q_j(r)$, uniformly in $z$, where $Q_j$ is the reference subtree law. For a parent started at $x$, the probability of exactly $l$ ordered child records $r_1,\ldots,r_l$ is
\begin{equation}\label{eq:record-kernel}
\int\left[\prod_{i=1}^l
U_j(\xi_{i-1},dz_i)S_j(z_i;r_i,d\xi_i)\right]
V_j(\xi_l,dy),\qquad \xi_0=x.
\end{equation}
For $l=0$, the integral is simply the total mass of $V_j(x,dy)$. In general, integrate first the final exit kernel, then the last child kernel, then the preceding entrance kernel, and continue backwards. Each bound is uniform in the position still to be integrated. With $\delta_j=H_0/\rho_{j+1}$, this gives the reference weight
\[
2^{-(l+1)}\prod_{i=1}^lQ_j(r_i)
\]
within a factor $\exp\{\pm[C(l+1)\delta_j+\sum_i\eta(r_i)]\}$. This order of integration preserves the dependence between a child record and its exit position.

Summing the numbers $l+1$ over parents at level $j-1$ gives exactly $N_{j-1}+N_j$. Starting from the terminal comparison and iterating~\eqref{eq:record-kernel}, for observation radius $a<m/2$ the total logarithmic error is at most
\begin{equation}\label{eq:2.24}
\eta\le
C\sum_{j=k+1}^{N_R}(N_{j-1}+N_j)\frac{H_0}{\rho_{j+1}}
+CN_{N_R}\frac{a+1}{m}\log\frac ma.
\end{equation}
This compares the record inside a root uniformly over its entrance; repeated application gives the corresponding comparison for a fixed forest. The reference forest starts with $N_k$ individuals and has independent marks of law~\eqref{eq:2.14} at its terminal generation. The exterior kernels determining the actual number of roots are treated in the next subsection.

If the generation bounds make the right side of~\eqref{eq:2.24} at most a deterministic $\eta_*$, summing the pointwise comparison over any set of permitted records preserves the factor $e^{\pm\eta_*}$. The set may impose joint conditions on all marks and generation counts. The same statement holds for every nonnegative functional supported on these records, by summation and monotone convergence on the countable record space.

\subsection{Returns to the root and a single-root lower bound}\label{sec:roots}
The forest comparison is uniform over root entrances. To sum over the number of roots, we combine it with the exterior kernels that determine returns after each root exit. Let $I(dx)$ be the subprobability kernel of the first root entrance, $S(x,dr,d\xi)$ the joint kernel of one root record and its exit, and $R(\xi,dx)$ the kernel of returning to the root entrance ball before $\tau_R$. Put $q(\xi)=1-R\mathbf1(\xi)$. The joint record law with exactly $j\ge1$ roots is
\begin{equation}\label{eq:root-record-kernel}
\begin{split}
&\mathbb P_0(J=j,dr_1\cdots dr_j)\\
&\quad=\int I(dx_1)
\left[\prod_{i=1}^{j-1}S(x_i,dr_i,d\xi_i)R(\xi_i,dx_{i+1})\right]
S(x_j,dr_j,d\xi_j)q(\xi_j).
\end{split}
\end{equation}
If $R\mathbf1\le\rho<1$, integrate from the final factor backwards. The first-entrance and final-nonreturn masses are at most one, each intervening return has mass at most $\rho$, and the root comparisons are uniform in the next entrance. Consequently, for every record event $E$ under the stated count bounds,
\begin{equation}\label{eq:2.25}
\mathbb P_0(J=j,E)\le \rho^{j-1}e^\eta\mathbf P_j^{\rm GW,mark}(E).
\end{equation}
The fixed shift in the return bound can be chosen uniformly over root levels. Write $H=H_0$, $a_*=A_{k+1}=\rho_{k+1}-H$ and $c_*=B_k=\rho_k+H$. For a root center $b\in D_R$ and an exit point still inside $D_R$, comparison with $B(b,2R)$ and the potential-kernel stopping argument give, once $a_*$ exceeds a fixed threshold,
\[
\mathbb P_{\rm exit}(\tau_R<H_{B(b,a_*)})
\ge\frac{\log(c_*/a_*)-4/a_*}{\log(2R/a_*)}.
\]
Here the inverse-square remainder in~\eqref{eq:2.5} can be bounded by $g/a_*$ beyond a fixed radius. The unit boundary overshoot and this remainder contribute at most $3g/a_*$ to the stopped potential; including its error at the starting point gives the displayed numerator error $4/a_*$. This fixes the coefficient independently of the root level.
Since $H\ge8$, $H\le a_*$ and $c_*\ge e(a_*+H)$, the buffer absorbs that error:
\[
\log(c_*/a_*)\ge1+\log(1+H/a_*)
\ge1+\frac{H}{a_*+H}\ge1+\frac4{a_*}.
\]
Also $R\le32e\rho_0$ and $a_*\ge\rho_{k+1}/2$ give
\[
\log(2R/a_*)\le k+\log32+2\log2+2.
\]
Choosing a fixed constant $K\ge\log32+2\log2+3$ proves $R\mathbf1\le1-(k+K)^{-1}$, and hence~\eqref{eq:2.25} with this return factor. The radius threshold and $H\le a_*$ follow from the standing choice $m\gg H_0$ with $m$ large; no condition growing with the total depth is needed.

For the lower bound, use root centers in $R/4\le|b|\le R/3$. Here $\rho_0\in[R/(32e),R/32]$, the starting point is outside the root, and the outer ball lies strictly in $D_R$. Estimate the first entrance in a disk centered at $b$ that contains the origin and lies inside $D_R$, and estimate nonreturn using a disk that contains $D_R$. Equation~\eqref{eq:2.10} gives $I\mathbf1\ge c_1$ and $\inf_\xi q(\xi)\ge c_2$, with fixed positive constants. Taking $j=1$ in~\eqref{eq:root-record-kernel} and using the lower record comparison therefore yields
\begin{equation}\label{eq:2.26}
c e^{-\eta}\mathbf P_1^{\rm GW,mark}(E)
\le\mathbb P_0(J=1,E)
\le e^\eta\mathbf P_1^{\rm GW,mark}(E).
\end{equation}
The event $J=1$ remains part of~\eqref{eq:2.26}, since its nonreturn factor $q(\xi)$ depends on the root exit. For pairs of sites we will need to prescribe exits as well; Section~\ref{sec:tworoot} supplies the corresponding bridge comparison.



\section{Branching-process and Gaussian estimates}\label{sec:tools}
The excursion comparison separates the spatial estimates into a branching-process barrier and a terminal occupation event. This section records the barrier and Gaussian estimates used for these two parts, and proves a comparison that bounds the terminal occupation event by a shifted Gaussian maximum.

\subsection{A geometric Galton--Watson barrier estimate}\label{sec:gw}
The reference offspring law from Section~\ref{sec:recordcomparison} is geometric. Let $Z_l$ be the critical Galton--Watson process with offspring distribution $\mathbb P(\xi=k)=2^{-(k+1)}$, and start it at $Z_0=x^2/2\in\mathbb N$. Write
\[
f_{a,b}(l;L)=a+(b-a)l/L,\qquad l_L=l\wedge(L-l).
\]
We use the following fixed-parameter form of~\cite[Theorem 9.1]{GW}. Fix $\delta>0$, $C\ge0$, $\eta>1$ and $\epsilon\in(0,1/2)$. For $\sqrt2\le x,y\le\eta L$, $x\le a$ and $y\le b$, the probability of staying below $f_{a,b}(l;L)+Cl_L^{1/2-\epsilon}$ and ending in $[y,y+\delta]\cap\sqrt{2\mathbb N}$ is at most
\begin{equation}\label{eq:3.1}
C_0\frac{(1+a-x)(1+b-y)}L
\sqrt{\frac{x}{yL}}\exp\left[-\frac{(x-y)^2}{2L}\right].
\end{equation}
Here and below the constants may depend on the fixed parameters. There is a lower bound of the same form, with the square-root factor replaced by its minimum with one, for paths constrained to the tube
\begin{equation}\label{eq:3.2}
[f_{x,y}-\widetilde C l_L^{1/2+\epsilon},\quad
 f_{a,b}-Cl_L^{1/2-\epsilon}],
\end{equation}
provided $(1+a-x)(1+b-y)\le\eta L$, $\max(xy,|y-x|)\ge L/\eta$, and the positive-lattice nonemptiness condition in the cited theorem holds. We check these requirements for the applications in Sections~\ref{sec:uppergw} and~\ref{sec:coarseonepoint}. The truncation of the square-root factor is inactive there. The source theorem has a closed terminal interval; when a half-open window is needed for a lower bound, we use a closed subinterval strictly inside it.

\subsection{Gaussian maxima and convex-set normal approximation}\label{sec:gaussianinputs}
A discrete Gaussian free field (DGFF) in a finite domain is the centered Gaussian field whose covariance is its SRW Green function. After division by $\sqrt g$, the maximum in a square of side $r$, with $h=\log r$, has centering
\[
m_h=2h-\tfrac34\log h+O(1).
\]
We use~\cite[Lemma 3.8, equation (39)]{BDZ} in the form
\begin{equation}\label{eq:3.3}
\mathbb P(\max\phi/\sqrt g\ge m_h+z)
\le C(1+z)e^{-2z-cz^2/h},\qquad z\ge1.
\end{equation}


To apply~\eqref{eq:3.3} at general side lengths, we compare with a dyadic square. Embed the smaller square in a dyadic square of at most a fixed multiple of its side. The larger field restricted to the smaller square is the smaller zero-boundary field plus an independent harmonic field. At a measurably chosen maximizer of the smaller field, the harmonic field is nonnegative with conditional probability at least $1/2$. The comparison therefore costs a fixed factor and a bounded change in centering.

We also use the quadratic increment estimate for the harmonic field in~\cite[Lemma 3.10, equation (45)]{BDZ}. Its constants are uniform in the large-square size and the number of subboxes, with the prescribed interior buffer. We use this estimate after embedding the terminal disk in a dyadic square in Section~\ref{sec:terminalupperproof}.

For normal approximation,~\cite[Theorem 1.1]{BE} states that, for independent centered $p$-dimensional vectors with total covariance equal to the identity, the error over any convex set is at most
\begin{equation}\label{eq:3.4}
(42p^{1/4}+16)\sum_i\mathbb E|X_i|^3.
\end{equation}
In Section~\ref{sec:tilt}, we apply~\eqref{eq:3.4} after whitening the mark vectors and compare the resulting error with the terminal-window probability.

\subsection{Proof of Lemma~\ref{lem:oneway}}\label{sec:oneway}
The next lemma transfers the Gaussian upper tail to discrete occupation counts. The holding-time comparison is applied only at a maximizer of those counts, so its cost is independent of the size of the observation set.

\begin{lemma}\label{lem:oneway}
Let $H=\sum_{i=1}^kY_i$ for $k\ge1$ independent marks of law~\eqref{eq:2.14}, and let $\phi$ be a centered Gaussian field of covariance $G_D$. For every finite $A\subset K$ and integer $a\ge1$,
\begin{equation}\label{eq:3.5}
\mathbf P_k(\max_AH\ge a)
\le C\mathbb P\left(\max_A|\phi+\sqrt{2\alpha k}|\ge\sqrt{2a}\right).
\end{equation}
\end{lemma}
\begin{proof}
First replace $k$ by an independent Poisson variable $Q$ of mean $k$. Since marks are nonnegative, the exceedance event is increasing in their number. For a nondegenerate centered random variable $X$ with a finite fourth moment, H\"older's and Cauchy--Schwarz inequalities give
\[
\mathbb P(X>0)\ge
\frac{(\mathbb E X^2)^2}{4\mathbb E X^4}.
\]
Indeed, $\mathbb E X_+=\tfrac12\mathbb E|X|$, and
$\mathbb E|X|\ge(\mathbb E X^2)^{3/2}/(\mathbb E X^4)^{1/2}$.
The centered second and fourth moments of $Q$ are $k$ and $3k^2+k$, so $\mathbb P(Q\ge k)\ge1/16$. The fixed-number probability is therefore bounded by a constant times the Poissonized probability.

Attach an independent mean-one exponential holding time to each visit of the Poissonized collection. Conditional on the discrete paths, choose a maximizer $z_*$ of the visit count by a deterministic rule. If the count there is $v\ge1$, its holding-time sum has law $\Gamma(v,1)$. The centered second and fourth moments are $v$ and $3v^2+6v$, and the preceding inequality gives $\mathbb P(\Gamma(v,1)\ge v)\ge1/36$. Thus passing to continuous occupation costs only another fixed factor and does not change the threshold.

Let $\mathcal L$ be the resulting occupation vector, $u=\alpha k$, and $F_0=\operatorname{diag}(f_z)$ for $f_z\ge0$. Substituting $v_z=(1+f_z)^{-1}$ into~\eqref{eq:2.16} and Poissonizing gives
\begin{equation}\label{eq:3.6}
\mathbb E e^{-\langle f,\mathcal L\rangle}
=\exp\{-u\mathbf1^TF_0(I+GF_0)^{-1}\mathbf1\}.
\end{equation}
Gaussian integration gives
\[
\mathbb E\exp\{-\tfrac12\langle f,(\phi+\sqrt{2u})^2\rangle\}
=\det(I+GF_0)^{-1/2}
\exp\{-u\mathbf1^TF_0(I+GF_0)^{-1}\mathbf1\}.
\]
Thus, for an independent field $\phi_0$ of covariance $G$,
\begin{equation}\label{eq:3.7}
\mathcal L+\tfrac12\phi_0^2
\overset{\rm law}=\tfrac12(\phi+\sqrt{2u})^2.
\end{equation}
The added term on the left is coordinatewise nonnegative. Combining this with the two preceding comparisons proves~\eqref{eq:3.5}.
\end{proof}
For the lower bound, we instead keep the fixed number of discrete marks and change measure using their generating function from Section~\ref{sec:markformulas}.



\section{Proof of Proposition~\ref{prop:upper}}\label{sec:upper}
We prove the upper bound by separating high points according to whether their traversal counts cross a curved barrier. Above the barrier, a spatial sum of traversal tails is sufficient. Below it, the branching estimate and a Gaussian bound for each terminal block give the required probability after summing over root returns. We work first in the central half of the disk, then obtain the full-domain estimate by enlarging the killing disk.

\subsection{Choice of scales}\label{sec:upperscales}
Choose the terminal radius so that the barrier suppresses crossings at the last coarse scales, while the record-comparison error remains negligible. Throughout this section use
\begin{equation}\label{eq:4.1}
\begin{gathered}
t=\log R,\quad D_+=(8\log t)^4,\quad m=e^{D_+},\quad
r=\lfloor m^{1/4}\rfloor,\quad H_0=4r+10,\\
\ell=\left\lfloor\log\frac R{32m}\right\rfloor,\quad N_R=\ell-1,
\quad d=g^{-1}G_{B(0,em-1)}(0,0),\quad T=N_R+d,\quad h=\log r.
\end{gathered}
\end{equation}
By~\eqref{eq:2.7},
\begin{equation}\label{eq:4.2}
d=D_++1+\kappa/g+o(1),\quad T=t+\Delta_R,
\quad |\Delta_R|\le C,\quad h/d\to1/4.
\end{equation}
In particular,
\begin{equation}\label{eq:4.3}
d^{1/4}=8\log t+o(1),\quad d^2/t\to0,\quad d\log t/t\to0,
\quad t^2D_+e^{-3D_+/4}\to0,
\quad R^2e^{-c_0t^2/d}\to0,
\end{equation}
for every fixed $c_0>0$. The effective depth $T$ therefore differs from $t$ by a bounded amount, including the rounding of the coarse depth.

For $0\le s\le\log t$, set
\begin{equation}\label{eq:4.4}
y=\sqrt\pi(b_R+s)=2t-\log t+\sqrt\pi s,\quad
\theta=y/T,\quad x_*=\theta N_R,\quad k_R(s)=\lceil(b_R+s)^2\rceil.
\end{equation}
We use this range of $s$ in the intermediate estimates; the final bound takes $s$ fixed before letting $R\to\infty$. Write $\lambda=2\sqrt\pi$.

\subsection{A terminal-block Gaussian estimate}\label{sec:terminalupper}
The next estimate combines fluctuations across a terminal block with the common harmonic contribution from larger scales. We retain the dependence on $q$ because the required height varies with the number of terminal marks.
\begin{lemma}\label{lem:terminalupper}
Let $C$ be a square of side comparable to $r$ in a fixed proportional interior of the terminal disk, and let $\phi$ have the SRW Green covariance of that disk. Uniformly for $q\ge-\sqrt d$,
\begin{equation}\label{eq:4.5}
\mathbb P\left(\max_C\frac{\phi}{\sqrt g}\ge2d+q\right)
\le C\frac{r^2e^{-2d}}d
\left(1+\frac{q_+}{\sqrt d}\right)e^{-2q-c_0q_+^2/d}.
\end{equation}
\end{lemma}

\subsection{Proof of Lemma~\ref{lem:terminalupper}}\label{sec:terminalupperproof}
Embed the terminal disk in a dyadic square of side comparable to $m$. The larger DGFF is the original field plus an independent harmonic field on the disk. At a chosen maximizer of the original field, the added field is nonnegative with probability at least $1/2$, so it suffices to work in the larger square. Choose a dyadic square $Q$ of side comparable to $r$ that contains $C$ with a fixed proportional buffer. The enclosing dyadic square can be translated so that $Q$ is one of its dyadic tiling cells and the terminal disk stays in its proportional interior. On $C$, decompose the normalized large field as the independent zero-boundary field $F$ in $Q$ plus its harmonic complement $H$.

The fine-field centering is $m_h=2h-\tfrac34\log h+O(1)$. Let $V$ be the harmonic field at a central lattice point and $v=\Var V=d-h+O(1)$. The Green diagonal estimates and the harmonic increment bound in Section~\ref{sec:gaussianinputs} give
\[
\sup_{x\in C}|\Var H_x-v|\le C,
\qquad
\Var(H_x-H_y)\le C|x-y|^2/r^2.
\]
In particular,
\[
\Cov(H_x,V)-v
=\tfrac12\bigl(\Var H_x-v-\Var(H_x-V)\bigr)=O(1).
\]
Gaussian regression yields
\begin{equation}\label{eq:4.6}
H_x=a_xV+\eta_x,\qquad |a_x-1|\le C/v,
\end{equation}
where $\eta$ is independent of $V$. Its value at the central point is zero, and
$\Var(\eta_x-\eta_y)\le\Var(H_x-H_y)$.
To control $A_0=\max_C|\eta_x|$, use dyadic nets anchored at that point. At level $j$ there are $O(4^j)$ increments with standard deviations at most $C2^{-j}$. The Gaussian tail bound and a union bound, with thresholds
$C(\sqrt{j+1}+u)2^{-j}$, show that all these thresholds hold with probability at least $1-Ce^{-cu^2}$. Summing the thresholds gives
\[
\mathbb P(A_0>C(1+u))\le Ce^{-cu^2},\qquad u\ge0.
\]
For large $v$, all $a_x$ are positive. Thus, with $V'=(1+C/v)V$ for a sufficiently large fixed $C$,
\[
\max_C H\le V'_++A_0,\qquad \Var(V')=v'=v+O(1),
\]
and $V'$ remains independent of the residual field.

Set $B=(\max_CF-m_h)_++A_0$. The two summands are independent. On $A_0\le z$, the inequality
$(z-A_0)^2\ge z^2/2-A_0^2$ and~\eqref{eq:3.3} give
\[
\begin{split}
\mathbb P(B\ge z,\ A_0\le z)
&\le C(1+z)e^{-2z-cz^2/(2h)}
 \mathbb E e^{2A_0+cA_0^2/h}\\
&\le C(1+z)e^{-2z-cz^2/(2h)}.
\end{split}
\]
The tail bound for $A_0$ makes the last expectation uniformly bounded for large $h$. The contribution from $A_0>z$ has the same upper bound after decreasing the quadratic constant. Hence
\begin{equation}\label{eq:4.7}
\mathbb P(B\ge z)\le C(1+z)e^{-2z-c_1z^2/h},\qquad z\ge0.
\end{equation}

Let $a=2d+q-m_h=2v+q+\tfrac34\log h+O(1)$. Since $q\ge-\sqrt d$, we have $a\ge c_2d$. Conditioning on $B$, and integrating its tail by parts, bounds
$\mathbb P(V'_++B\ge a)$ by $\mathbb P(V'\ge a)$, $\mathbb P(B\ge a)$, and
\begin{equation}\label{eq:4.8}
\frac C{\sqrt{v'}}\int_0^a(1+z)
\exp\left[-\frac{(a-z)^2}{2v'}-2z-c_1z^2/h\right]dz.
\end{equation}
Completing the square separates the leading exponential cost:
\begin{equation}\label{eq:4.9}
-\frac{(a-z)^2}{2v'}-2z
=-2a+2v'-\frac{(a-2v'-z)^2}{2v'}.
\end{equation}
The resulting factor is
\[
C d^{-1/2}h^{-3/2}e^{-2(d-h)-2q}
\]
outside the remaining integral. If $-\sqrt d\le q\le0$, dropping the first nonpositive square bounds that integral by
$\int_0^\infty(1+z)e^{-c_1z^2/h}\,dz\le Cd$.
For $q>0$, write $\zeta=a-2v'=q+\tfrac34\log h+O(1)$ and use
\begin{equation}\label{eq:terminal-double-square}
\frac{(\zeta-z)^2}{2v'}+\frac{c_1z^2}{h}
=\frac{h+2c_1v'}{2v'h}
 \left(z-\frac{h\zeta}{h+2c_1v'}\right)^2
 +\frac{c_1\zeta^2}{h+2c_1v'}.
\end{equation}
Here $v'\asymp h\asymp d$. Integration of the Gaussian and its first absolute moment therefore gives
\[
\int_0^a(1+z)
 e^{-(\zeta-z)^2/(2v')-c_1z^2/h}\,dz
\le Cd(1+q/\sqrt d)e^{-c_3q^2/d}.
\]
Multiplying by the prefactor and using
$d^{1/2}h^{-3/2}\asymp d^{-1}$ and
$e^{-2(d-h)}=r^2e^{-2d}$, gives the right-hand side of~\eqref{eq:4.5}.

The Gaussian Mills bound and~\eqref{eq:4.9} also give
\[
\mathbb P(V'\ge a)
\le C d^{-1/2}e^{-2a+2v'}
 e^{-(a-2v')^2/(2v')}
\le C\frac{r^2e^{-2d}}{d^{1/2}h^{3/2}}
 e^{-2q-c_4q_+^2/d},
\]
which is bounded by the same right-hand side. Finally,
$\mathbb P(B\ge a)\le C(1+a)e^{-2a-c_1a^2/h}$.
For $q\le0$, use $a\ge c_2d$; for $q>0$, use $a\asymp d+q$. Relative to the required bound, the extra factor $e^{-2v'}$ absorbs the polynomial in $a$, while part of $e^{-c_1a^2/h}$ retains the decay in $q^2/d$. Decreasing $c_0$ if necessary proves~\eqref{eq:4.5} over its full stated range.
\qed

\subsection{Upper tails for terminal occupation marks}\label{sec:terminalmarksupper}
We now apply Lemma~\ref{lem:oneway} to the terminal marks. Recall that $N_{N_R}$ is their number in the reference forest, and put $u=\sqrt{2N_{N_R}}$. If $u\ge t-1$ and
\[
u\le x_*+A+d^{1/4},\qquad A=4+s,
\]
set $q=x_*-u$. Then $q\ge-A-d^{1/4}$, and this lower bound exceeds $-\sqrt d$ by a quantity tending to infinity, uniformly for $0\le s\le\log t$. Since $\alpha=g+O(r/m)$, $\sqrt{2k_R(s)/g}=y+O(t^{-1})$, and $y-x_*=\theta d$, the Gaussian threshold in Lemma~\ref{lem:oneway} satisfies
\begin{equation}\label{eq:4.10}
\begin{split}
\sqrt{2k_R(s)/g}-\sqrt{\alpha/g}\,u
&=2d+q+\varepsilon_R,\\
|\varepsilon_R|
&\le C\left(t^{-1}+tr/m+d\log t/t\right)=o(1).
\end{split}
\end{equation}
The error bound is uniform over the displayed range of $u,s$, including $q_+=O(t)$. The inequality
\[
(q+\varepsilon_R)_+^2\ge\tfrac12q_+^2-\varepsilon_R^2
\]
allows a decrease of the quadratic constant in~\eqref{eq:4.5}; its linear exponential and polynomial factors change by only a bounded factor. Applying Lemma~\ref{lem:terminalupper} to the positive branch therefore gives
\begin{equation}\label{eq:4.11}
\mathbf P_{N_{N_R}}(\max_CH\ge k_R(s))
\le C\frac{r^2e^{-2d}}d
\left(1+\frac{q_+}{\sqrt d}\right)e^{-2q-c_0q_+^2/d}
+Cr^2e^{-c_1t^2/d}.
\end{equation}
The last term bounds the negative branch of the absolute value in~\eqref{eq:3.5}. If $u<t-1$, the entire probability is at most $Cr^2e^{-c_1t^2/d}$, since the distance to the shifted Gaussian threshold is of order $t$. We can therefore sum over every terminal count allowed by the coarse barrier, using the Gaussian error alone for the smaller counts.

\subsection{Spatial sum of barrier violations}\label{sec:barrierviolations}
We next bound the probability that any point-centered traversal count crosses the barrier. Use
\begin{equation}\label{eq:4.12}
B_{j,R}(s)=\theta j+A+w_j,\qquad
w_j=(j\wedge(T-j))^{1/4},\quad 1\le j\le N_R,
\end{equation}
where $A=4+s$. In particular, its curvature at the last coarse generation is $w_{N_R}=d^{1/4}$.

For each $j$, construct the mesh
\[
a_j=\left\lfloor\frac{\rho_{j+1}}{1000(j+1)}\right\rfloor,
\quad \Lambda_j=a_j\mathbb Z^2\cap B(0,3R/4),
\quad h_j=2a_j+3.
\]
At a mesh point $z$, use inner radius $b_j=\rho_{j+1}+h_j$ and outer radius $c_j=\rho_j-h_j$. Define $\widetilde N_{j,z}$ by the entrance--exit recursion~\eqref{eq:traversal-stopping}, with center $z$ and radii $b_j,c_j$, counting every inner entrance before $\tau_R$. The first entrance is at time zero if the origin is in the inner ball. Each counted segment is complete: its outer ball, with a one-step buffer, lies in $D_R$.

Assign $x\in B(0,R/2)$ to its nearest mesh point $z$, so that $|x-z|\le a_j$. The choice of $h_j$ gives
\[
B(x,\rho_{j+1})\subset B(z,b_j),
\qquad B(z,c_j)\subset B(x,\rho_j).
\]
Assign each inward record about $x$ to the mesh segment containing its inner-entrance time. The first inclusion makes the assignment well defined; the second ensures that an outer-shell visit between two records separates their mesh segments. Thus
\begin{equation}\label{eq:4.13}
\mathcal T_j^x\le\widetilde N_{j,z}.
\end{equation}

Since $D_R\subset B(z,2R)$, equation~\eqref{eq:2.10} bounds the escape probability after an outer mesh exit from below. Uniformly in the exit point, its numerator is at least
$1-\epsilon/(j+1)$ for a fixed $\epsilon<1/4$, and its denominator is at most $j+C_0$. Here the coefficient $1000$ makes the geometric error small, while the lattice error is smaller still because $m/t\to\infty$. Choosing a fixed $K$ sufficiently large gives
\[
\mathbb P_{\rm exit}(\tau_R<H_{B(z,b_j)})
\ge\frac{1-\epsilon/(j+1)}{j+C_0}
\ge\frac1{j+K}.
\]

For the first entrance, if $|z|>2b_j$, apply~\eqref{eq:2.10} in the same containing disk, starting at the origin:
\[
\mathbb P_0(H_{B(z,b_j)}<\tau_R)
\le C\frac{1+\log(2R/|z|)}{1+\log(2R/b_j)}.
\]
If $|z|\le2b_j$, the trivial bound by one suffices. Since
$\log(2R/b_j)=j+O(1)$, both cases give
\[
\mathbb P_0(H_{B(z,b_j)}<\tau_R)
\le\frac{C\psi_j(z)}{j+1},
\qquad
\psi_j(z)=1+\log\frac{2R}{|z|\vee b_j}.
\]
For $u>0$, the strong Markov property at successive outer exits now bounds the tail by the first-entrance probability times
$(1-(j+K)^{-1})^{\lceil u^2/2\rceil-1}$. Consequently,
\begin{equation}\label{eq:4.14}
\mathbb P_0(\sqrt{2\widetilde N_{j,z}}\ge u)
\le\frac{C\psi_j(z)}{j+1}
\exp\left[-\frac{u^2}{2(j+K)}\right],
\qquad u>0.
\end{equation}

To sum the hitting weights, note that $b_j/a_j\asymp j+1$ and
$R/a_j\asymp(j+1)e^j$. The mesh points with $|z|\le2b_j$ contribute at most $C(j+1)^3$. In a dyadic annulus of radius $2^{-\nu}R>2b_j$, there are at most $C(2^{-\nu}R/a_j)^2$ points and $\psi_j(z)\le C(\nu+1)$. Hence
\begin{equation}\label{eq:4.15}
\begin{aligned}
\sum_{z\in\Lambda_j}\psi_j(z)
&\le C(j+1)^3+C(R/a_j)^2\sum_{\nu\ge0}(\nu+1)4^{-\nu}\\
&\le C(j+1)^2e^{2j}.
\end{aligned}
\end{equation}

For $v\in\mathbb N_0$, let
$E_{j,z,v}=\{\sqrt{2\widetilde N_{j,z}}\ge B_{j,R}(s)+v\}$.
We use the identity
\begin{equation}\label{eq:4.16}
\frac{(\theta j+U)^2}{2(j+K)}
=\frac{\theta^2j}{2}+\theta U-\frac{\theta^2K}{2}
+\frac{(U-\theta K)^2}{2(j+K)}.
\end{equation}
Together with $\theta\ge3/2$ and
\[
e^{(2-\theta^2/2)j}\le C(j+1)^2,\qquad 1\le j\le N_R,
\]
this gives
\begin{equation}\label{eq:4.17}
\sum_{z\in\Lambda_j}\mathbb P_0(E_{j,z,v})
\le C(j+1)^3e^{-3(A+w_j+v)/2}.
\end{equation}
For $j\le N_R/3$, the sum is controlled by $\sum_j(j+1)^3e^{-cj^{1/4}}<\infty$. The middle range contributes at most $Ct^4e^{-ct^{1/4}}$. In the final range, $w_j\ge d^{1/4}$, and the total is at most
\[
Ct^4e^{-3d^{1/4}/2}=t^{-8+o(1)}.
\]
The sum over $v$ is geometric. Define
\[
\mathcal B_R(s)=\{\exists x\in B(0,R/2),\ \exists j\le N_R:
\sqrt{2\mathcal T_j^x}>B_{j,R}(s)\}.
\]
The preceding sum gives
\begin{equation}\label{eq:4.18}
\limsup_{R\to\infty}\mathbb P_0(\mathcal B_R(s))\le Ce^{-3s/2}.
\end{equation}
Thus~\eqref{eq:4.18} controls all barrier crossings, regardless of the local time at the crossing point.

\subsection{Root levels determined by the distance to the origin}\label{sec:rootlevels}
For a block near the origin, the walk may start inside the largest excursion shells. We choose its root level so that the entrance ball excludes the origin, which allows the exact decomposition~\eqref{eq:2.21} to apply. Cover $R/t<|x|\le R/2$ by half-open lattice squares of side $r$. For a block center $b_C$, choose the smallest $k\ge0$ such that
\begin{equation}\label{eq:4.19}
2\rho_{k+1}+H_0\le|b_C|.
\end{equation}
Then $0\le k\le\log t+C$. The block family $\mathscr C_{R,k}$ with this root level obeys
\begin{equation}\label{eq:4.20}
|\mathscr C_{R,k}|\le C R^2e^{-2k}/r^2.
\end{equation}
The origin is outside the root entrance ball and all later observed shells. Thus the pathwise decomposition~\eqref{eq:2.21} and the count comparison~\eqref{eq:2.22} apply for $j\ge k+1$, even if the origin lies inside the root outer ball.

For the root returns, the quantitative estimate in Section~\ref{sec:roots} applies with the same fixed $K$: the containing disk is $B(b_C,2R)$, $H_0\ge10$, and $A_{k+1}\ge m-H_0\to\infty$ uniformly over the root levels under consideration. Thus the return bound holds with
\begin{equation}\label{eq:4.21}
\varrho_k=1-(k+K)^{-1}.
\end{equation}
On $J_C\le t^2$ and $N_j^C\le16t^2$,~\eqref{eq:2.24} gives
\begin{equation}\label{eq:4.22}
\eta_R\le Ct^2\left(\frac rm+\frac{r+1}m\log\frac mr\right)
\le Ct^2D_+e^{-3D_+/4}=o(1).
\end{equation}
The strong Markov property also gives
$\mathbb P_0(J_C>t^2)\le\varrho_k^{\lfloor t^2\rfloor}$.
Summing~\eqref{eq:4.20} over the root levels bounds the total number of blocks by $CR^2/r^2\le CR^2$. Therefore their total probability of exceeding the root cap is at most
\begin{equation}\label{eq:4.23}
CR^2e^{-ct^2/\log t}=o(1),
\end{equation}
because $2t-ct^2/\log t\to-\infty$.
Below the barrier, all other generation counts are at most $16t^2$. Define
\begin{equation}\label{eq:4.24}
W_C^+(s)=\{\max_{z\in C}L_R(z)\ge k_R(s),\quad
\sqrt{2N_j^C}\le B_{j,R}(s),\ k+1\le j\le N_R\}.
\end{equation}
By~\eqref{eq:2.21}, the local time in this event is the sum of the terminal marks from all $J_C$ roots.

\subsection{The branching barrier and the root-number sum}\label{sec:uppergw}
For a block with root level $k$ and $J$ roots, the reference forest has depth $N_R-k$ and initial population $J$. Let $L=N_R-k$, $x=\sqrt{2J}$, and consider a terminal window $[u,u+1)$, where $u\ge t-1$ and $q=x_*-u$. In~\eqref{eq:3.1}, use
\begin{equation}\label{eq:4.25}
a=x+A+\theta k+2k^{1/4},\quad
b=x_*+A+2d^{1/4},\quad C=1,\ \epsilon=1/4,\ \eta=4,\ \delta=1.
\end{equation}
The source barrier dominates the actual barrier. In the first half, use
$(k+l)^{1/4}\le k^{1/4}+l^{1/4}$; in the second half, use
$(L-l+d)^{1/4}\le(L-l)^{1/4}+d^{1/4}$. The remaining linear endpoint margins are nonnegative. The conditions $J\le t^2$, $L\sim t$, and $u\sim t$ place the states in the source range. Also $b-u=A+2d^{1/4}+q\ge0$. Enlarging the half-open endpoint window to the closed source window, we obtain
\begin{equation}\label{eq:4.26}
\mathbf P_J(\text{barrier},\sqrt{2Z_L}\in[u,u+1))
\le\frac{C(1+A+k)S(q)}{L^2}J^{1/4}
 e^{-u^2/(2L)+(u/L)\sqrt{2J}},
\end{equation}
where $S(q)=1+A+2d^{1/4}+q$.

We now sum this estimate against the geometric return bound~\eqref{eq:4.21}. For $q_0\ge1$ and $a\in[1/2,3]$,
\begin{equation}\label{eq:4.27}
\sum_{J\ge1}(1-q_0^{-1})^{J-1}J^{1/4}e^{a\sqrt{2J}}
\le Cq_0^2e^{a^2q_0/2}.
\end{equation}
Indeed, $(1-q_0^{-1})^{J-1}\le e\,e^{-J/q_0}$, with the term $q_0=1,J=1$ interpreted as one, and
\[
-J/q_0+a\sqrt{2J}
=-\frac{(\sqrt J-aq_0/\sqrt2)^2}{q_0}+a^2q_0/2.
\]
With $v=\sqrt J$, the sum is bounded by a constant times
\[
1+\int_0^\infty (v+1)^{3/2}
 e^{-(v-aq_0/\sqrt2)^2/q_0}\,dv\le Cq_0^2.
\]
Take $q_0=k+K$ and $a=u/L\le2+C\log t/t$. Since $k=O(\log t)$,
\begin{equation}\label{eq:4.28}
\sum_{J\ge1}\varrho_k^{J-1}J^{1/4}e^{(u/L)\sqrt{2J}}
\le C(k+1)^2e^{2k}.
\end{equation}
The factor $e^{2k}$ is canceled by the block-count factor $e^{-2k}$ in~\eqref{eq:4.20}.

\subsection{Summation over terminal windows and root levels}\label{sec:uppersum}
The terminal-mark bound~\eqref{eq:4.11}, the branching-window bound~\eqref{eq:4.26}, and the root sum~\eqref{eq:4.28} now control the event $W_C^+(s)$. Apply the forest comparison~\eqref{eq:4.22} and sum over the blocks in~\eqref{eq:4.20}. After separating the absolute errors, which are estimated below, the contribution from root level $k$ is at most
\begin{equation}\label{eq:4.29}
\begin{split}
&C(1+A+k)(k+1)^2\frac{e^{2t-2d}}{(N_R-k)^2}\\
&\quad\times\sum_{q\ge-Q}\frac{S(q)}d
(1+q_+/\sqrt d)
\exp\left[-\frac{(x_*-q)^2}{2(N_R-k)}-2q-c_0q_+^2/d\right],
\end{split}
\end{equation}
Here the sum is over unit terminal windows, and $Q=A+d^{1/4}+3$. Rounding their endpoints changes only a fixed factor.

Writing $T=t+\Delta_R$ gives
\begin{equation}\label{eq:4.30}
\frac{x_*^2}{2N_R}
=2t-2d-2\log t+\lambda s-2\Delta_R
+O\left(\frac{d\log t}{t}+\frac{(\log t)^2}{t}\right).
\end{equation}
Thus $e^{2t-2d-x_*^2/(2N_R)}/N_R^2\le Ce^{-\lambda s}$. The reduced forest depth at root level $k$ contributes the additional exponential cost
\begin{equation}\label{eq:4.31}
\frac{x_*^2}{2(N_R-k)}-\frac{x_*^2}{2N_R}
=\frac{\theta^2 N_R k}{2(N_R-k)}\ge9k/8.
\end{equation}
The remaining coefficient of $q$ is
$\beta_k=x_* /(N_R-k)-2=O(\log t/t)$, and $d\beta_k^2=o(1)$, so part of the quadratic decay in $q/\sqrt d$ absorbs it.

Set $S_0=1+A+2d^{1/4}+3$. Summing separately over the positive and negative values of $q$ gives
\begin{equation}\label{eq:4.32}
\frac1d\sum_{q\ge-Q}(S_0+q)(1+q_+/\sqrt d)
 e^{\beta_kq-c_1q_+^2/d}
\le C\left[1+\frac{S_0}{\sqrt d}+\frac{(Q+1)S_0}{d}\right]\le C.
\end{equation}
The last bound is uniform for $0\le s\le\log t$, because $Q,S_0=O(\log t)$ and $\sqrt d\asymp(\log t)^2$. Hence the sum of the principal contributions is at most
\begin{equation}\label{eq:4.33}
Ce^{-\lambda s}\sum_{k\ge0}(1+A+k)(k+1)^2e^{-9k/8}
\le C(1+s)e^{-\lambda s}.
\end{equation}

For the low-terminal-state and negative-Gaussian-branch error, first integrate $Cr^2e^{-c_0t^2/d}$ against the entire reference endpoint distribution. Then sum over roots, using $\sum_J\varrho_k^{J-1}\le k+K$, and finally over blocks. The resulting error is at most $CR^2e^{-c_1t^2/d}=o(1)$. Together with~\eqref{eq:4.23}, this proves
\begin{equation}\label{eq:4.34}
\limsup_{R\to\infty}\sum_C\mathbb P_0(W_C^+(s))
\le C(1+s)e^{-\lambda s}.
\end{equation}

\subsection{Proof of Proposition~\ref{prop:upper}}\label{sec:uppercompletion}
Estimate~\eqref{eq:4.34} controls high points away from the origin whose traversal counts stay below the barrier. We combine it with the barrier-crossing bound~\eqref{eq:4.18}, a one-point estimate near the origin, and an enlargement of the disk to include its boundary region. By~\eqref{eq:2.7}--\eqref{eq:2.9}, for $|x|\le R/t$,
\[
\mathbb P_0(L_R(x)\ge k_R(s))
\le CR^{-2}t e^{-\lambda s}
\left(1+\log\frac R{|x|\vee1}\right).
\]
Summing radially and treating the origin separately gives, for each fixed $s$,
\begin{equation}\label{eq:4.35}
\mathbb P_0\left(\max_{|x|\le R/t}L_R(x)\ge k_R(s)\right)
\le Ce^{-\lambda s}\frac{1+\log t}{t}+o(1)=o(1).
\end{equation}
Every high point in $B(0,R/2)$ either lies in this small neighborhood, violates one of the traversal barriers, or belongs to a block satisfying~\eqref{eq:4.24}, by~\eqref{eq:2.22}. Thus
\begin{equation}\label{eq:4.36}
\mathbf1_{\{\max_{|x|\le R/2}L_R(x)\ge k_R(s)\}}
\le\mathbf1_{\{\text{a high point near the origin}\}}+\mathbf1_{\mathcal B_R(s)}
+\sum_C\mathbf1_{W_C^+(s)}.
\end{equation}
Equations~\eqref{eq:4.18},~\eqref{eq:4.34} and~\eqref{eq:4.35} imply
\begin{equation}\label{eq:4.37}
\limsup_{R\to\infty}\mathbb P_0\left(\max_{|x|\le R/2}\sqrt{L_R(x)}\ge b_R+s\right)
\le C\{(1+s)e^{-\lambda s}+e^{-3s/2}\}.
\end{equation}
Recall that $\mathfrak c=2/\sqrt\pi$. On the same infinite walk,
\[
\begin{gathered}
M_R\le\max_{|x|\le R}L_{2R}(x),\\
b_{2R}-b_R
=\mathfrak c\log2-\frac1{\sqrt\pi}
 \log\left(1+\frac{\log2}{\log R}\right)
<\mathfrak c\log2.
\end{gathered}
\]
For fixed $s\ge\mathfrak c\log2$, this gives the event inclusion
\[
\{\sqrt{M_R}>b_R+s\}
\subset
\left\{\max_{|x|\le R}\sqrt{L_{2R}(x)}
\ge b_{2R}+s-\mathfrak c\log2\right\}.
\]
Apply~\eqref{eq:4.37} at radius $2R$ with the fixed height
$s-\mathfrak c\log2$. The bounded shift changes only the constant in the claimed tail bound. Increasing that constant also covers
$0\le s<\mathfrak c\log2$, and proves Proposition~\ref{prop:upper}.
\qed



\section{Proof of Proposition~\ref{prop:onepoint}}\label{sec:events}
We construct events on which a specified site has local time at least $b_R^2$, and prove that each event has probability of order $R^{-2}$. Besides a barrier on the excursion counts, the event imposes a continuous constraint on equilibrium averages around the site. The latter constraint will control the shared part of the occupation profiles when we estimate nearby pairs in Section~\ref{sec:second}.

\subsection{Construction of successful sites}\label{sec:lowerscales}
We first specify the scales separating the excursion counts from the terminal occupation record. For the remainder of the paper, set
\begin{equation}\label{eq:5.1}
\begin{gathered}
t=\log R,\quad D_-=(2000\log t)^{5/2},\quad m=e^{D_-},\quad r=1,\quad H_0=14,\\
\ell=\left\lfloor\log\frac R{32m}\right\rfloor,\quad N_R=\ell-1,
\quad\rho_j=me^{\ell-j},\quad
d=g^{-1}G_{B(0,em-1)}(0,0),\quad T=N_R+d.
\end{gathered}
\end{equation}
Thus $d=D_-+O(1)$, $T=t+O(1)$, $N_R\sim t$, and
\begin{equation}\label{eq:5.2}
d^2/t+d\log t/t\longrightarrow0.
\end{equation}
Put $\gamma=2/5$ and
\begin{equation}\label{eq:5.3}
y=2t-\log t=\sqrt\pi b_R,\quad \theta=y/T,\quad
x_*=\theta N_R,\quad k_R=\lceil b_R^2\rceil,\quad
w_j=(j\wedge(T-j))^\gamma.
\end{equation}
In particular $2-\theta=(\log t+O(1))/t>0$ and
$w_{N_R}=d^\gamma=2000\log t+o(1)$.
Let $\mathscr Z_R=\{z\in\mathbb Z^2:R/4\le |z|\le R/3\}$.
At each $z\in\mathscr Z_R$, let $J_z$ and $N_j^z$ be the root and generation counts from
Section~\ref{sec:physical}, with the concentric disks centered at $z$. Define the coarse barrier event by
\begin{equation}\label{eq:5.4}
\mathcal G_z=\{J_z=1;\ \sqrt{2N_j^z}\le4+\theta j-w_j,
\quad 1\le j\le N_R\}.
\end{equation}
For $k=N_{N_R}^z>0$, write $u=\sqrt{2k}$ and $q=x_*-u$. We restrict the deficit $q$ at the last generation to the band
\begin{equation}\label{eq:5.5}
q\in[\sqrt d,2\sqrt d].
\end{equation}
All generation counts on $\mathcal G_z$ are at most $Ct^2$.

\label{sec:averageconstraints}
Inside the last excursion scale, we observe the occupation at $z$ together with averages on nested equilibrium supports. These averages will have the covariance of Brownian motion up to a small diagonal correction. The terminal domain and reference intensity are
\begin{equation}\label{eq:5.6}
\Omega_z=B(z,em+H_0),\qquad
\alpha=g\log\frac{em+H_0}{m-H_0}=g+O(m^{-1}).
\end{equation}
Fix $B>4$ sufficiently large, and then fix a sufficiently large $A$.
For $0\le i\le I$ put
\begin{equation}\label{eq:5.7}
a_i=AB^i,\qquad I=\max\{i:a_i\le m/t^6\},\qquad a_{\max}=a_I,
\end{equation}
and, with all capacities relative to $\Omega_z$, put
\begin{equation}\label{eq:5.8}
\omega_i=\frac{e_{B(z,a_i)}}{\operatorname{cap}(B(z,a_i))},\quad
v_i=\frac1{g\operatorname{cap}(B(z,a_i))}-\frac\alpha{2g},\quad
\varepsilon_i=\frac{\|\omega_i\|_2^2}{2g}.
\end{equation}
The equilibrium supports are disjoint and avoid $z$. Equations~\eqref{eq:2.11}--\eqref{eq:2.12} give
\begin{equation}\label{eq:5.9}
v_i=d-\log a_i+O(1),\quad \varepsilon_i\le C/a_i,\quad
v_I=6\log t+O(1),\quad I=O(d).
\end{equation}
Choosing $B$ first bounds each difference $v_i-v_{i+1}$ above and below by positive fixed constants.
For every observed site $x$, center and normalize the original visit counts by setting
\[
F_x=\frac{L_R(x)-\alpha k}{u\sqrt{\alpha g}},\qquad
S_i=\omega_i^TF,\qquad
V=\frac{G_{\Omega_z}(z,z)}g-\frac{1+\alpha}{2g}.
\]
The remaining variance between the innermost average and the site will stay of constant order. More precisely, the deterministic quantities $H_*=d-v_0=\log A+O(1)$ and
$\Delta=V-v_0=\log A+O(1)$ are positive for $A$ large. Set $a_*=\sqrt{H_*}$.
The constants implicit here are independent of $A$; after fixing $A$, both $H_*$ and $\Delta$ remain in fixed compact subintervals of $(0,\infty)$.

\label{sec:success}
To impose the constraint between observation times, enlarge the probability space by independent standard Brownian bridges, one on each interval between consecutive times in $0<v_I<\cdots<v_0$. The bridges for different sites are independent of one another and of the walk. Let $\mathcal B_z$ be the linear interpolation of
$0,S_I-2v_I,\ldots,S_0-2v_0$, plus these bridges. Define $\mathcal T_z(k)$ by
\begin{equation}\label{eq:5.10}
\begin{gathered}
A_t(k)\le F_z\le A_t(k)+1,\qquad
A_t(k)=\frac{k_R-\alpha k}{u\sqrt{\alpha g}},\\
q-a_*-1\le S_0-2v_0\le q-a_*,\\
\mathcal B_z(v)\le q-(d-v)^\gamma\qquad(0\le v\le v_0).
\end{gathered}
\end{equation}
In particular $S_i\le2v_i+q-(d-v_i)^\gamma$ at every observation time.
Set
\begin{equation}\label{eq:5.15}
\mathcal C_z=\mathcal G_z\cap\{q\in[\sqrt d,2\sqrt d]\}\cap\mathcal T_z(k),
\qquad W_R=\sum_{z\in\mathscr Z_R}\mathbf1_{\mathcal C_z},
\end{equation}
and set $\mathcal C_z=\varnothing$ when $k=0$. The auxiliary variables only restrict the event: pathwise,
\begin{equation}\label{eq:5.16}
\{W_R>0\}\subset\{M_R\ge k_R\}\subset\{\sqrt{M_R}\ge b_R\}.
\end{equation}
All probabilities of success below include the auxiliary variables. The next proposition gives the first moment required for the second-moment argument; its proof occupies the rest of this section.

\begin{proposition}\label{prop:onepoint}
For all sufficiently large real $R$, uniformly over $z\in\mathscr Z_R$,
\begin{equation}\label{eq:7.28}
cR^{-2}\le\mathbb P_0(\mathcal C_z)\le CR^{-2},\qquad
0<c_1\le\mathbb EW_R\le C_1.
\end{equation}
\end{proposition}

\label{sec:lowercomparison}
We will calculate the terminal event under the reference mark law. To justify this change of law, note that all observations lie inside $B(z,a_{\max})\subset B(z,m-H_0-2)$, so the complete-record identity~\eqref{eq:2.21} includes every visit contributing to them.
Under the coarse count bounds the comparison error in~\eqref{eq:2.24} is
\begin{equation}\label{eq:5.17}
\eta_R^{\rm term}\le Ct^2/m+
Ct^2\frac{a_{\max}+1}{m}\log\frac m{a_{\max}}
\le Ct^2/m+Ct^{-4}\log t=o(1).
\end{equation}
The positive-mark mass is $\asymp1/\log t<1/2$, so~\eqref{eq:2.26}
applies to the full event. The auxiliary bridges are included by first fixing their paths, applying the uniform record comparison, and then integrating over their common independent law. We use the same procedure for the conditional record comparisons in the pair estimates.

\subsection{Proof of Lemma~\ref{lem:gaussian-density}}\label{sec:covariance}
The probability of the terminal event tends to zero, so we will compare Gaussian probabilities by a density ratio. We first identify the two covariance matrices entering that comparison.
Under the law $\mathbf P_k$ of $k=u^2/2$ independent reference marks, use
the same notation $F,S_i$ with $L_R$ replaced by their summed occupations.
The exact identity~\eqref{eq:2.19} and the constant equilibrium potential give
\begin{equation}\label{eq:single-covariance}
\begin{gathered}
\Var S_i=v_i-\varepsilon_i,\quad
\Cov(S_i,S_j)=\min(v_i,v_j)\ (i\ne j),\\
\Cov(F_z,S_i)=v_i,\qquad \Var F_z=V=d+O(1).
\end{gathered}
\end{equation}
Indeed, if the support of one average lies inside the ball defining the other, the latter equilibrium potential is constant on that support and equals the reciprocal capacity. The discrete diagonal correction vanishes on disjoint supports. Thus, in the order
$\mathsf F=(S_I,\ldots,S_0,F_z)$, the covariance is exactly
\begin{equation}\label{eq:single-KE}
\Sigma=K-E,\qquad K_{ij}=s_i\wedge s_j,\qquad
(s_1,\ldots,s_p)=(v_I,\ldots,v_0,V),\qquad
E=\diag(\varepsilon_I,\ldots,\varepsilon_0,0).
\end{equation}
The quadratic form of the inverse Brownian covariance is
\[
x^TK^{-1}x=\frac{x_1^2}{s_1}
+\sum_{j=2}^p\frac{(x_j-x_{j-1})^2}{s_j-s_{j-1}}.
\]
All denominators are bounded below uniformly, so $K\succeq cI$. Since
$\|E\|\le C/A$, enlarging $A$ gives $\Sigma\succeq c'I$, uniformly in $p=O(d)$.

\begin{lemma}\label{lem:gaussian-density}
For the matrices in~\eqref{eq:single-KE}, let
$J=E^{1/2}K^{-1}E^{1/2}$. The centered Gaussian densities satisfy
\begin{equation}\label{eq:density-ratio}
\frac{f_\Sigma(x)}{f_K(x)}=\det(I-J)^{-1/2}e^{-Q(x)/2},\qquad
0\le Q(x)\le2\sum_i\varepsilon_i((K^{-1}x)_i)^2,
\end{equation}
and $1\le\det(I-J)^{-1/2}\le e^{C/A}$.
For any event $\mathcal E$, possibly including independent auxiliary variables,
$\mathbb P_\Sigma(\mathcal E)\le e^{C/A}\mathbb P_K(\mathcal E)$.
If $\mathbb E_K[Q\mid\mathcal E]\le C_A$, then also
$\mathbb P_\Sigma(\mathcal E)\ge e^{-C_A/2}\mathbb P_K(\mathcal E)$.
\end{lemma}
\begin{proof}
The matrix inverse identity gives
$Q=x^TK^{-1}E^{1/2}(I-J)^{-1}E^{1/2}K^{-1}x$.
Here $\|J\|\le1/2$ for $A$ large and
$\operatorname{tr}J\le C\sum_i\varepsilon_i\le C/A$.
The determinant identity gives~\eqref{eq:density-ratio}. The pointwise upper bound yields the probability upper bound, while conditional Jensen yields the lower bound. Independent auxiliary variables can be included because their law is unchanged.
\end{proof}

\subsection{Proof of Lemma~\ref{lem:bounded-tilt}}\label{sec:tilt}\label{sec:tiltnormal}
Exponential tilting brings the large terminal occupation into the range of a Gaussian approximation. We establish an estimate that also permits the signed tilt needed for pairs. Suppose $\mathsf F$ consists of $p\le Cd$ sites or probability-weighted averages of $F$, with covariance $\Sigma\succeq cI$. If $\beta\in\mathbb R^p$ satisfies $\|\beta\|_1\le C$, put
\[
\Lambda_\beta=\log\mathbf E_k e^{\beta^T\mathsf F},\qquad
\frac{d\mathbf P_k^\beta}{d\mathbf P_k}
=e^{\beta^T\mathsf F-\Lambda_\beta}.
\]
\begin{lemma}\label{lem:bounded-tilt}
Uniformly for $u\asymp t$ and the observations just described,
\begin{equation}\label{eq:6.12}
\Lambda_\beta=\tfrac12\beta^T\Sigma\beta+O(d^3/t),
\end{equation}
and, for every convex Borel set $D\subset\mathbb R^p$,
\begin{equation}\label{eq:6.13}
\left|\mathbf P_k^\beta(\mathsf F\in D)
-\mathbb P(N(\Sigma\beta,\Sigma)\in D)\right|
\le\delta_t,\qquad \delta_t=Cd^5/t=o(d^{-a})
\end{equation}
for each fixed $a>0$. The estimate remains valid after fixing independent auxiliary variables when the resulting section of the event is convex, and then integrating them.
\end{lemma}
\begin{proof}
Let $U$ be the observation vector of one mark and $c_u=u\sqrt{\alpha g}$.
The tilt is a product of single-mark densities proportional to
$\exp(\beta^TU/c_u)$, so independence is preserved. The geometric law~\eqref{eq:2.18}, Jensen and H\"older give
$\mathbb E U_i^6\le Cd^5$ and
$\mathbb E\exp(bU_i)\le1+Cb$ for $0\le b\le c/d$.
These bounds apply also to $\sum_i|\beta_i|U_i/\|\beta\|_1$.
Cauchy--Schwarz, together with Jensen's inequality for the normalizing denominator, therefore gives a uniform bound $Cd^3$ for every third centered coordinate moment under all intermediate single-mark tilts. H\"older's inequality gives the same bound for each third joint cumulant. Taylor expansion of the logarithmic moment generating function proves~\eqref{eq:6.12}. Differentiating once and twice gives, entrywise,
\begin{equation}\label{eq:tilt-cumulant-errors}
|\mu_i-(\Sigma\beta)_i|+|\Gamma_{ij}-\Sigma_{ij}|
\le Cd^3/t,
\end{equation}
where $\mu,\Gamma$ are the tilted mean and covariance of $\mathsf F$.
Since $\|\Gamma-\Sigma\|\le Cpd^3/t=o(1)$, $\Gamma\succeq cI/2$.
After whitening the centered independent summands, their Lyapunov sum is at most
$Ckp^{3/2}d^3/c_u^3\le Cp^{3/2}d^3/t$.
Equation~\eqref{eq:3.4} bounds the convex-set error by
$Cp^{7/4}d^3/t$.
Finally~\eqref{eq:tilt-cumulant-errors} bounds the standardized mean error by
$C\sqrt p\,d^3/t$ and the covariance error in Frobenius norm by $Cpd^3/t$.
The Gaussian relative entropy formula and Pinsker's inequality give a total variation error $Cpd^3/t$. Combining this with the convex-set error gives the bound $Cd^5/t$.
Conditioning and integrating independent auxiliary variables preserves this uniform bound.
\end{proof}

\subsection{Proof of Lemma~\ref{lem:bridge}}\label{sec:bridge}
The Gaussian calculation reduces the constraint on the averages to a Brownian bridge that remains below a curve. The endpoint stays a fixed distance from the flat barrier; this produces a factor of order $d^{-1/2}$ in the terminal probability.
\begin{lemma}\label{lem:bridge}
Let $v_0=d-H_*$, where $H_*=\log A+O(1)$, and let $a_*=\sqrt{H_*}$.
For $A$ fixed sufficiently large, all large $d$, $q\in[\sqrt d,2\sqrt d]$, and
$b\in[q-a_*-1,q-a_*]$, a standard Brownian bridge from $0$ to $b$ in time $v_0$ satisfies
\begin{equation}\label{eq:7.12}
\mathbb P(B(v)\le q-(d-v)^\gamma\ (0\le v\le v_0)\mid B(v_0)=b)
\asymp\frac{qa_*}{d}.
\end{equation}
Under this conditioned law its increments have second moments at most
$C_A(s-r)$ on every $[r,s]\subset[0,v_0]$.
\end{lemma}
\begin{proof}
Writing $a=q-b\in[a_*,a_*+1]$, reflection gives the flat-barrier probability
\begin{equation}\label{eq:7.14}
1-e^{-2qa/v_0}\asymp qa_*/d.
\end{equation}
Conditioned on this flat event, $X(u)=q-B(v_0-u)$ is a BES$(3)$ bridge from $a$ to $q$ of duration $v_0$. Indeed the killed Brownian and BES$(3)$ densities are
\begin{equation}\label{eq:7.15}
p_s^+(x,y)=\frac{e^{-(x-y)^2/(2s)}-e^{-(x+y)^2/(2s)}}{\sqrt{2\pi s}},
\qquad q_s(x,y)=\frac yx p_s^+(x,y),
\end{equation}
and the factors $y/x$ cancel in bridge densities.
The explicit formula implies $q_s(x,y)\le Cy^2s^{-3/2}$ for all $x,y,s>0$.
Since $aq/v_0\to0$ and $(q-a)^2/v_0$ is bounded, it also implies
$q_{v_0}(a,q)\asymp q^2v_0^{-3/2}$ and
$q_{v_0}(q,a)\asymp a^2v_0^{-3/2}$.
Consequently, on its first half the bridge has density at most a universal constant relative to BES$(3)$ from $a$. Reversibility,
$x^2q_s(x,y)=y^2q_s(y,x)$, gives the same domination for the reversed second half relative to BES$(3)$ from $q$.

For an unconditioned BES$(3)$ process started at $x>0$, write $H_b=\inf\{u\ge0:X_u=b\}$ for $b>0$. Then
\begin{equation}\label{eq:7.17}
\mathbb P_x(H_b<\infty)=\min(1,b/x),\qquad
\mathbb E_x X_s^{-1}
=x^{-1}\mathbb P_x(\tau_0^{\rm BM}>s)\le C/\sqrt s.
\end{equation}
Here is a finite-time derivation of the hitting identity. For $0<b<x$ and $T>0$, let $s_T$ be the survival mass up to time $T$ of Brownian motion started at $x-b$ and killed at zero. The relation $q_T(x,y)=(y/x)p_T^+(x,y)$ extends to the path law by cancellation of the intermediate factors. Restricting the Brownian path to stay above $b$ and shifting by $b$, the killed density has mass $s_T$ and first moment $x-b$. Therefore
\[
\mathbb P_x\!\left(\min_{0\le u\le T}X_u\ge b\right)
=\frac{x-b+b s_T}{x},\qquad
\mathbb P_x\!\left(\min_{0\le u\le T}X_u<b\right)
=\frac bx(1-s_T).
\]
Reflection gives $s_T=\mathbb P(|N(0,T)|\le x-b)\to0$. Increasing the time horizon gives infinite-time strict lower hitting with probability $b/x$; continuity and squeezing nearby levels give the same probability for hitting the level $b$. Integrating $q_s(x,y)/y$ in~\eqref{eq:7.15} gives the Brownian survival identity in~\eqref{eq:7.17}, and reflection bounds it by $C/\sqrt s$. Finally, if $x<b$, then $\mathbb P_x(X_T\le b)\le b\mathbb E_xX_T^{-1}\to0$, so continuity implies that $b$ is reached almost surely; if $x=b$, it is reached at time zero. This proves both identities in~\eqref{eq:7.17} with all starting positions covered.
The curved restriction becomes $X(u)\ge(H_*+u)^\gamma$.
On $[0,H_*]$ its failure has probability at most $CH_*^\gamma/a\le CH_*^{\gamma-1/2}$.
On $[L,2L]$, $L=2^jH_*$, condition at $L$ and use~\eqref{eq:7.17}; the probability is at most $CL^\gamma\mathbb E_aX_L^{-1}\le CL^{\gamma-1/2}$.
Summing up to $v_0/2$ gives $CH_*^{\gamma-1/2}$.
On the other half the boundary is at most $d^\gamma$; viewed backwards from $q$, its hitting probability is at most $d^\gamma/q\le d^{\gamma-1/2}$.
Using the density bounds on the two halves, the probability of crossing the curve, conditional on staying below the flat barrier, is at most
$C(H_*^{\gamma-1/2}+d^{\gamma-1/2})$.
Choose $A$ so that the first term is below $1/4$, and then $d$ so that the second is below $1/4$. This proves~\eqref{eq:7.12}.

An unconditioned BES$(3)$ process is $|x e_1+W^{(3)}|$; hence
$\mathbb E(X_s-X_r)^2\le3(s-r)$.
The half-path dominations give this bound up to a constant under the flat bridge; an interval crossing the midpoint is split in two and the square inequality is used. The curved restriction has conditional probability at least $1/2$, so its further conditioning only changes the constant. Time reversal gives the asserted increment bound for $B$.
\end{proof}

\subsection{Proof of Proposition~\ref{prop:terminalprob}}\label{sec:onepoint}\label{sec:terminalprob}
We now combine the tilted normal approximation with the bridge estimate. The endpoint window and the bridge constraint each cost a factor of order $d^{-1/2}$; the inverse tilt supplies the exponential term.
\begin{proposition}\label{prop:terminalprob}
Uniformly over the integer values $k=u^2/2$ with $q=x_*-u\in[\sqrt d,2\sqrt d]$,
\begin{equation}\label{eq:7.18}
\mathbf P_k(\mathcal T_z(k))\asymp d^{-1}e^{-2d-2q}.
\end{equation}
\end{proposition}
\begin{proof}
The exact threshold satisfies
\begin{equation}\label{eq:7.19}
A_t(k)=2d+q+o(1).
\end{equation}
Indeed $(y^2-u^2)/(2u)=y-u+(y-u)^2/(2u)$, with error $O(d^2/t)$; also
$y-u=\theta d+q$, $\theta d-2d=O(d\log t/t)$, and the intensity and rounding errors are $O(t/m+t^{-1})$.
Under the reference mark law $\mathbf P_k$, tilt by $2F_z$. Lemma~\ref{lem:bounded-tilt} gives
$\Lambda=2V+O(d^3/t)$, Gaussian mean
$(2v_I,\ldots,2v_0,2V)$, covariance $\Sigma$, and error $\delta_t=o(d^{-1})$.
For fixed auxiliary bridges, the restrictions~\eqref{eq:5.10} form a convex set: the continuous inequality is an intersection of affine half-spaces. Thus the lemma applies to this event.

For the Gaussian comparison, subtract the mean $(2v_I,\ldots,2v_0,2V)$ and first use covariance $K$. The interpolated average coordinates then form a Brownian motion $B$, and
$B(v_0)\in[q-a_*-1,q-a_*]$ has probability $\asymp d^{-1/2}$.
Lemma~\ref{lem:bridge} supplies a conditional factor $\asymp d^{-1/2}$.
The centered site coordinate is $B(v_0)$ plus an independent increment of variance
$\Delta=V-v_0$. The site window in~\eqref{eq:5.10} requires that increment to lie in a unit interval with endpoints $a_*+O(1)$. Since $\Delta$ and $a_*$ stay in fixed compact ranges depending only on $A$, its probability is bounded above and below by positive constants. The Gaussian event therefore has probability $\asymp d^{-1}$.

It remains to check the conditional moment required for the lower bound in Lemma~\ref{lem:gaussian-density}. Let $X$ denote the centered Gaussian vector with covariance $K$. On the Brownian event just described, the final-increment restriction changes the distribution of $B(v_0)$ within its window by a bounded factor. Lemma~\ref{lem:bridge} consequently bounds the second moment of every earlier increment by $C_A$ times its duration. The final increment is bounded in absolute value by $C_A$.
Each component of $K^{-1}X$ is the difference of at most two adjacent increments divided by their time lengths (with the first term $X_1/s_1$). All those lengths are bounded below. Hence
\[
\sup_i\mathbb E_K[((K^{-1}X)_i)^2\mid\mathcal T_z(k)]\le C_A,
\qquad \mathbb E_K[Q\mid\mathcal T_z(k)]\le C_A\sum_i\varepsilon_i\le C_A.
\]
Lemma~\ref{lem:gaussian-density} now gives probability $\asymp d^{-1}$ also for covariance $\Sigma$.
The convex-set error is $o(d^{-1})$, so the terminal event also has probability $\asymp d^{-1}$ under the tilted reference mark law. On the unit site window the inverse tilt is
$\exp(\Lambda-2F_z)\asymp e^{-2d-2q}$, proving~\eqref{eq:7.18}.
\end{proof}

\subsection{Proof of Proposition~\ref{prop:onepoint}}\label{sec:coarseonepoint}\label{sec:firstmomentcompletion}
We return to the full success event $\mathcal C_z$, which combines the coarse barrier with the terminal occupation constraint. For a geometric Galton--Watson process started from one individual, and $u=x_*-q$ in the terminal band, the estimates of Section~\ref{sec:gw} give
\begin{equation}\label{eq:7.24}
\mathbf P_1\bigl(\sqrt{2Z_j}\le4+\theta j-w_j\ (1\le j\le N_R),
\ \sqrt{2Z_{N_R}}\in[u,u+1)\bigr)
\asymp\frac{1+q}{N_R^2}e^{-u^2/(2N_R)}.
\end{equation}
To verify the hypotheses of the lower tube estimate, choose
$L=N_R$, $x=\sqrt2$, source endpoint $y_s=u+1/4$,
$a=4$, $b=4+\theta N_R-2d^\gamma$, $C=1$, $\widetilde C=10$,
$\epsilon=1/10$, $\eta=4$, and $\delta=1/2$.
The closed endpoint window $[u+1/4,u+3/4]$ lies inside the desired half-open window.
The upper curve in that estimate lies below the barrier defining $\mathcal G_z$. In the first half, its linear part is decreased; in the second half, use
$(j\wedge(T-j))^\gamma\le(N_R-j)^\gamma+d^\gamma$ and
$2d^\gamma j/N_R\ge d^\gamma$.
The endpoint margins have product $O(q)=o(N_R)$ and final margin $\asymp q$, because $d^\gamma=o(\sqrt d)$. The product $xy_s$ is of order $N_R$, and both states satisfy the bounds required by the tube estimate.
The two endpoint lines are ordered, and at every interior layer the difference of their curvature terms is at least $9$. The upper curve is at least $4+\tfrac32j-j^{2/5}\ge4$ for large $R$, so the tube contains a positive state of $\sqrt{2\mathbb N}$; adjacent states are at distance at most $\sqrt2$.
At the terminal height $u\asymp t$, their spacing is $O(t^{-1})$, so the closed endpoint window also contains a state. This verifies the tube and lattice hypotheses. The Gaussian exponential at the shifted endpoint differs by only a bounded factor from that in~\eqref{eq:7.24}.
For the upper bound apply~\eqref{eq:3.1} to endpoint barrier heights $4,4+x_*$, with positive curvature and a closed unit window containing the desired window.

To combine the coarse and terminal probabilities, expand the coarse exponential as
\begin{equation}\label{eq:7.26}
\frac{x_*^2}{2N_R}=2t-2d-2\log t+O(1),\qquad
\frac{(x_*-q)^2}{2N_R}=\frac{x_*^2}{2N_R}-\theta q+\frac{q^2}{2N_R}.
\end{equation}
After multiplication by the terminal exponential in~\eqref{eq:7.18}, the remaining factor $e^{(\theta-2)q-q^2/(2N_R)}$ tends uniformly to one.
Multiply~\eqref{eq:7.24} by~\eqref{eq:7.18} and sum over unit terminal windows. The probabilities vary within each window by at most a fixed factor, and
\begin{equation}\label{eq:7.27}
d^{-1}\sum_{q\in[\sqrt d,2\sqrt d]\text{ in unit windows}}(1+q)\asymp1.
\end{equation}
The single-root comparison~\eqref{eq:2.26} with~\eqref{eq:5.17} gives
$\mathbb P_0(\mathcal C_z)\asymp R^{-2}$. Since $|\mathscr Z_R|\asymp R^2$, this proves Proposition~\ref{prop:onepoint}.\qed

\section{Pairs of sites that separate at an early scale}\label{sec:tworoot}
To bound the second moment of the successful-site count, we first consider pairs whose excursion records separate at an early generation. Conditional on a suitable exterior record, we compare their later records with two reference forests. This comparison leaves the two root counts under the joint law of the original walk; we estimate that law through the common part of the excursion record.

\subsection{Proof of Lemma~\ref{lem:continuation}}\label{sec:continuation}
Set
\begin{equation}\label{eq:7.29}
\mathcal L_t=\lceil(1000\log t)^{5/2}\rceil.
\end{equation}
For an integer $0\le\nu\le2\mathcal L_t$, start a geometric Galton--Watson forest with $k$ individuals at level $\nu$. Impose the coarse barriers in~\eqref{eq:5.4} after level $\nu$, the terminal band~\eqref{eq:5.5}, and the terminal mark event $\mathcal T_z(k_{N_R})$ from Section~\ref{sec:terminalprob}. Here $k_{N_R}$ is the terminal population, which fixes the number of marks in that event. Write $p_\nu(k)$ for the resulting reference probability, with the auxiliary bridges integrated out, and set $p_\nu(0)=0$.

\begin{lemma}\label{lem:continuation}
If $x=\sqrt{2k}\le4+\theta\nu$, then
\begin{equation}\label{eq:7.30}
p_{\nu}(k)\le C(\nu+1)^{3/2}\frac{1}{R^2}
 e^{-2\nu+2\sqrt{2k}}.
\end{equation}
\end{lemma}
\begin{proof}
The case $k=0$ is immediate. For $k\ge1$, apply the upper barrier estimate over length $L=N_R-\nu$, with endpoint heights $4+\theta\nu$ and $4+\theta N_R$, after enlarging the barrier by a fixed positive curvature. The product of the endpoint margins is at most $C(\nu+1)(1+q)$, and the state range in the estimate is valid because $\nu=O(\mathcal L_t)=o(t)$. Combining~\eqref{eq:3.1} with the terminal probability bound~\eqref{eq:7.18} and summing over the terminal band gives
\begin{equation}\label{eq:7.31}
p_{\nu}(k)\le C(\nu+1)\sqrt{x}\,\frac{e^{-2d}}{t^2}
\exp\left[-\frac{(x_*-x)^2}{2(N_R-\nu)}\right].
\end{equation}
The remaining linear coefficient of $q$ is $O((\mathcal L_t+\log t)/t)$. Its product with $\sqrt d$ tends to zero, so it is uniformly negligible on the terminal band. The window sum~\eqref{eq:7.27} then cancels the factor $d^{-1}$.

We use the identity
\begin{equation}\label{eq:7.32}
\frac{(\theta N_R-x)^2}{2(N_R-\nu)}
=\frac{\theta^2N_R}{2}+\frac{\theta^2\nu}{2}-\theta x
+\frac{(\theta \nu-x)^2}{2(N_R-\nu)}.
\end{equation}
Together with~\eqref{eq:7.26}, $\theta\le2$, $(2-\theta)\nu=o(1)$, and $\sqrt x\le C(\nu+1)^{1/2}$, it proves the result.
\end{proof}
To apply this reference estimate to two sites of the walk, we next compare their later records after conditioning on the exterior path.

\subsection{Removing the conditioning on the long exits}\label{sec:exterior}\label{sec:exitkernel}
The exterior path specifies the exits of certain long excursions. We remove this conditioning by comparing the complete traces inside smaller balls; the comparison error will be uniform in the number of visits. First, let a killed finite Markov chain in $D$ start at $x\in K\subset D$. Put
\[
 h_e(x)=\mathbb P_x(S_{\tau_D}=e),\quad
 q(x)=\mathbb P_x(\tau_D<H_K^+),\quad
 q_e(x)=\mathbb P_x(\tau_D<H_K^+,S_{\tau_D}=e).
\]
For every positive-probability trace $\xi=(x_0,\ldots,x_n)$ of all visits to $K$,
\begin{equation}\label{eq:trace-rn}
 \frac{\mathbb P_x(\xi\mid S_{\tau_D}=e)}{\mathbb P_x(\xi)}
 =\frac{q_e(x_n)}{h_e(x)q(x_n)}.
\end{equation}
Indeed, let $Q_K$ be the kernel of the next return to $K$. The trace has probability $\prod_{l<n}Q_K(x_l,x_{l+1})q(x_n)$; its joint probability with exit at $e$ is the same product with $q_e(x_n)$ in place of $q(x_n)$. Dividing these expressions and then dividing by $h_e(x)$ proves~\eqref{eq:trace-rn} for every trace length.

For SRW, if $D=B(b,\mathfrak b)$, $K\subset B(b,\mathfrak a)$ and $\mathfrak b/\mathfrak a$ is sufficiently large, then uniformly in reachable exit atoms $e$,
\begin{equation}\label{eq:trace-disk}
 \frac{q_e(x_n)}{h_e(x)q(x_n)}=1+O(\eta),\qquad
 \eta=\frac{\mathfrak a}{\mathfrak b}\left(1+\log\frac{\mathfrak b}{\mathfrak a}\right).
\end{equation}
The bound is relative to the exit probability, so it also applies to exit atoms of small probability. To prove it, put $h_*=h_e(b)$. The harmonic estimate~\eqref{eq:2.13} gives $h_e=h_*(1+O(\mathfrak a/\mathfrak b))$ on $B(b,3\mathfrak a)$. For $v\notin K$, put $Q(v)=\mathbb P_v(\tau_D<H_K)$ and $U(v)=\mathbb P_v(\tau_D<H_K,S_{\tau_D}=e)$. On the boundary of $B(b,2\mathfrak a)$, the identity
\[
 U(v)=h_e(v)-\mathbb E_v[h_e(S_{H_K});H_K<\tau_D]
\]
gives $|U-h_*Q|\le C(\mathfrak a/\mathfrak b)h_*$. The annulus escape estimate also gives $Q\ge c/(1+\log(\mathfrak b/\mathfrak a))$ there. Since both functions vanish on $K$, the maximum principle yields $U=h_*Q(1+O(\eta))$ throughout $B(b,2\mathfrak a)\setminus K$. Averaging over the first step from $K$ proves the corresponding bound for $q_e/q$; division by $h_e(x)$ proves~\eqref{eq:trace-disk}. The resulting measure inequalities remain valid for any function of the trace and after adjoining independent auxiliary variables.

For distinct candidate sites $z,w$, take short inner radius $A_{\nu+1}=\rho_{\nu+1}-H_0$, short outer radius $B_\nu=\rho_\nu+H_0$, and long outer radius $B_{\nu-J}$. We will choose the scales so that the two long balls have a one-step gap, lie in $D_R$, and avoid the origin. Wait for entrance into either short inner ball, delete the segment up to exit from the corresponding long ball, and repeat chronologically until killing. Let $\mathsf H$ record the exterior path words and the label, entrance and exit of each deleted segment, and put $\mathscr H=\sigma(\mathsf H)$. The record includes neither the deleted lengths nor absolute times; in particular, it does not determine $\tau_R$.

Conditional on an exterior value $h$ with $m$ deleted segments, their interiors are independent SRW bridges to the prescribed long exits. Indeed, the fiber over $h$ is the product of the sets of finite internal path words with the prescribed first exits. Inserting words $u_1,\ldots,u_m$ gives path weight $4^{-L_{\rm ext}}\prod_i4^{-|u_i|}$. Summing over all internal lengths gives
\begin{equation}\label{eq:exterior-fiber}
 \mathbb P_0((u_i)_i\mid\mathsf H=h)
 =\prod_{i=1}^m\frac{4^{-|u_i|}}{\mathbb P_{x_i}(S_{\tau_{D_i}}=e_i)}.
\end{equation}
The independent auxiliary bridges assigned to $z,w$ are also independent of this exterior record.

In each long segment, retain the trace on the ball of radius $B_\nu+2$. It determines every short root and its full fine record. The enlargement includes the one-step overshoot at each short exit. By~\eqref{eq:trace-disk}, removing the specified long exit costs $C(J+1)e^{-J}$ in the logarithm of the density per segment, uniformly in the number of short roots. For an unconditioned long segment, the escape probability following a short exit is
\begin{equation}\label{eq:8.1}
 p=\frac{\log(B_\nu/A_{\nu+1})}{\log(B_{\nu-J}/A_{\nu+1})}
 \asymp(J+1)^{-1},\qquad
 \mathbb P_\xi(\tau_{B_{\nu-J}}<H_{A_{\nu+1}})
 =p\{1+O((H_0+1)/A_{\nu+1})\}.
\end{equation}
Integrating the escape kernel, the last short-root kernel and the preceding return kernel backwards, as in~\eqref{eq:record-kernel}, gives the reference measure for $k\ge1$ ordered fine records $r_1,\ldots,r_k$:
\begin{equation}\label{eq:8.4}
 p(1-p)^{k-1}\prod_{i=1}^k\mathbf Q_\nu(dr_i)
\end{equation}
within logarithmic error
\[
 C(J+1)e^{-J}+Ck(H_0+1)/A_{\nu+1}+\sum_i\eta(r_i).
\]
Here $\mathbf Q_\nu$ is the reference tree below one level-$\nu$ root and $\eta(r_i)$ is the error in~\eqref{eq:2.24}. The kernel bounds are uniform in every intermediate entrance and exit, and each short-root record remains coupled to its exit until the later kernels have been integrated. Integrating instead over all fine records gives the same comparison for the number of roots, with the last error sum omitted.

\subsection{Keeping the original joint root-count law}\label{sec:jointroots}
The preceding comparison separates the fine records conditional on their root totals. We now integrate over those totals while retaining their original dependence. Let $K_z,K_w$ be the original total short-root counts and
$K_\nu^{\max}=\lfloor(4+\theta\nu)^2/2\rfloor\le C(\nu+1)^2$.
Let $E_z,E_w$ be the events that the remaining records succeed, including their independent auxiliary bridges. Given $\mathsf H=h$, write $M_z,M_w$ for the numbers of long segments. Under~\eqref{eq:8.4}, the total for label $E$ has law $\mu_{M_E}$, the law of a sum of $M_E$ geometric variables with parameter $p$; set $\mu_0=\delta_0$. Conditional on a total $k$, the fine records are $k$ independent copies of $\mathbf Q_\nu$, and their success probability, with the auxiliary bridges integrated out, is $p_\nu(k)$.

For $k,l\le K_\nu^{\max}$, apply the upper comparison for records to $E_z,E_w$, and the lower comparison for root counts after integrating over all fine records. There are at most $k+l$ long segments. The success events bound the total population at each fine generation, so~\eqref{eq:5.17} controls the sum of the fine-record errors. With
\begin{equation}\label{eq:8.5}
 \eta_{\rm dec}\le C(\nu+1)^2
 \{(J+1)e^{-J}+(H_0+1)/A_{\nu+1}\}+2\eta_R^{\rm term},
\end{equation}
these two comparisons respectively give
\[
 \mathbb P_0(E_z,E_w,K_z=k,K_w=l\mid h)
 \le e^{C\eta_{\rm dec}}\mu_{M_z}(k)\mu_{M_w}(l)p_\nu(k)p_\nu(l)
\]
and $\mu_{M_z}(k)\mu_{M_w}(l)\le e^{C\eta_{\rm dec}}\mathbb P_0(K_z=k,K_w=l\mid h)$.
Consequently, for every $G\in\sigma(\mathsf H)$,
\begin{equation}\label{eq:8.6}
\begin{split}
 &\mathbb P_0(G,E_z,E_w,K_z,K_w\le K_\nu^{\max})\\
 &\quad\le e^{C\eta_{\rm dec}}\mathbb E_0\!\left[
 \mathbf1_Gp_\nu(K_z)p_\nu(K_w)
 \mathbf1_{\{K_z,K_w\le K_\nu^{\max}\}}\right].
\end{split}
\end{equation}
The comparisons hold with the auxiliary bridge paths fixed and hence after integration against their independent laws. Monotone convergence justifies the sums over the countable exterior records and root totals. Thus the expectation in~\eqref{eq:8.6} is taken under the original joint law of $K_z,K_w$.

\subsection{Distance scales and an exterior-measurable prefix}\label{sec:earlygeometry}
For distinct sites, put $\rho=|z-w|$ and
\begin{equation}\label{eq:8.7}
a(z,w)=\left\lfloor\log\frac{\rho_0}{\rho}\right\rfloor.
\end{equation}
The early-distance range that varies with $t$ is $a<\mathcal L_t+J_t$, where
$J_t=\lceil10\log t\rceil+5$. Fix large constants $a_\circ,J_0$. For
$a_\circ<a<\mathcal L_t+J_t$, choose
\begin{equation}\label{eq:8.8}
J=\lceil20\log(a+2)\rceil+J_0,\qquad
j=a-J,\qquad \nu=a+J+5.
\end{equation}
By increasing $a_\circ$, we may assume $j\ge a/2>0$ and $\nu\le2\mathcal L_t$. The long outer radius is
$B_{\nu-J}=\rho_{a+5}+H_0\le e^{-4}\rho+o(\rho)$.
Thus the two long balls are disjoint, away from the origin, and contained in $D_R$.

Both long balls lie strictly inside the common entrance ball $B(z,\rho_{j+1}-H_0)$, whose radius is of order $e^J\rho$. This allows us to use the barrier through level $j$ as the exterior-measurable event in~\eqref{eq:8.6}. Define
\begin{equation}\label{eq:8.9}
G_j=\{J_z=1;\ \sqrt{2N_l^z}\le4+\theta l-w_l,\quad1\le l\le j\}.
\end{equation}
Every deleted long segment lies strictly inside all boundaries used by $G_j$. Removing or refilling it preserves the ordered boundary hits that determine $G_j$. Since these hits are recorded without absolute times, $G_j\in\mathscr H$.

On $\mathcal C_z\cap\mathcal C_w$, the short-root totals satisfy
$K_z=N_\nu^z$ and $K_w=N_\nu^w$: grouping by long excursions and the nested decomposition enumerate the same complete short segments. Equation~\eqref{eq:8.5} gives
\begin{equation}\label{eq:8.10}
\eta_{\rm dec}\le C(a+1)^2(J+1)e^{-J}+o(1)\le C.
\end{equation}
Using~\eqref{eq:8.6} and Lemma~\ref{lem:continuation}, and then dropping the root-count bounds on the right, gives
\begin{equation}\label{eq:8.11}
\begin{split}
\mathbb P_0(\mathcal C_z\cap\mathcal C_w)
\le C(\nu+1)^3\frac{1}{R^4}e^{-4\nu}
\mathbb E_0[\mathbf1_{G_j}\exp\{2\sqrt{2K_z}+2\sqrt{2K_w}\}],
\end{split}
\end{equation}
We may also impose $N_j^z\ge1$ on the right, since success is impossible when $N_j^z=0$. It remains to estimate the exponential moment of the two root counts under their original joint law.

\subsection{The common reference marks}\label{sec:commonmarks}
To estimate this joint moment, use the common record through level $j$. Observe a ball about $z$ whose radius is a fixed multiple of $\rho$ and which contains both long balls. Its ratio to the common entrance radius is $O(e^{-J})$. The barrier in $G_j$ bounds the number of marks at level $j$ by $O(j^2)$. Applying~\eqref{eq:2.24} to the common record therefore gives total comparison error
\begin{equation}\label{eq:8.12}
Cj^2\left(\frac{H_0}{\rho_{j+1}}+(J+1)e^{-J}\right)\le C.
\end{equation}
The probability of a positive mark is $O(1/J)<1/2$. Under the reference law, a common count $k$ gives $k$ independent marks, each retaining all crossings between the two short-root neighborhoods.

The common mark is killed in $\Omega_j=B(z,\rho_j+H_0)$ and has constant mean visit intensity
\begin{equation}\label{eq:8.13}
\alpha_j=g\log\frac{\rho_j+H_0}{\rho_{j+1}-H_0}=g+O(H_0/\rho_{j+1}).
\end{equation}
Both short inner and outer balls lie inside the observation ball with sufficient separation from the killing boundary.

\subsection{The first moment of the two short-root counts}\label{sec:rootmean}
For one common mark, let $K_E$ count the complete short roots of label $E=z,w$. Each short root has a unique last visit to its inner ball before exit from the corresponding short outer ball. Summing over ordinary visit times, applying the Markov property to their futures, and using the constant visit intensity $\alpha_j$ gives
\begin{equation}\label{eq:8.14}
\begin{split}
\mathbf EK_E
&=\sum_{x\in B(E,A_{\nu+1})}\mathbf E[\text{visits of the mark to }x\text{}]
\mathbb P_x(\tau_{B(E,B_{\nu})}<H_{B(E,A_{\nu+1})}^+)\\
&=\alpha_j\operatorname{cap}_{B(E,B_{\nu})}(B(E,A_{\nu+1})).
\end{split}
\end{equation}
This counts last visits by testing the future at each ordinary visit time. The mark starts outside both short balls, and every short root that starts is completed before exit from the common killing domain.

At the center of the inner ball,~\eqref{eq:2.4} and Green symmetry give, uniformly on the support of its equilibrium measure,
\[
G_{B(E,B_\nu)}(E,x)
=G_{B(E,B_\nu)}(x,E)
=g\log(B_\nu/A_{\nu+1})+O(A_{\nu+1}^{-1}).
\]
Therefore,
\begin{equation}\label{eq:8.15}
\operatorname{cap}_{B(E,B_{\nu})}(B(E,A_{\nu+1}))
=\{g\log(B_{\nu}/A_{\nu+1})+O(A_{\nu+1}^{-1})\}^{-1}.
\end{equation}
Since $\log(B_\nu/A_{\nu+1})=1+O(H_0/A_{\nu+1})$,
\begin{equation}\label{eq:8.16}
\mu:=\mathbf E(K_z+K_w)
=2+O(H_0/A_{\nu+1}+H_0/\rho_{j+1}).
\end{equation}
All physical radii in the early-distance range are at least $cRe^{-3\mathcal L_t}$. The error in~\eqref{eq:8.16} therefore remains negligible after multiplication by any fixed power of $\mathcal L_t$ used below. Keeping the leading coefficient $2$ will give the required exponential bound after summation over pairs.

\subsection{Factorial moments at completion times}\label{sec:factorial}
The short outer balls are disjoint, so roots of the two labels occur in chronological order without interlacing. Their completion times are stopping times at which no root remains unfinished. Let $\mathcal S$ be the set of possible exit positions from either short outer ball. Every point of $\mathcal S$ lies in $\Omega_j$ and outside both short outer balls. For $y\in\mathcal S$, let $K_y$ count all subsequent complete short roots of either label before exit from $\Omega_j$.

Write $I_E=B(E,A_{\nu+1})$ and $O_E=B(E,B_\nu)$ for $E=z,w$. The visit-sum identity used in~\eqref{eq:8.14}, now started at $y$, gives
\[
\mathbb E_yK_y
=\sum_{E=z,w}\sum_{x\in I_E}
G_{\Omega_j}(y,x)e_{I_E}^{O_E}(x).
\]
If $y$ is an exit of $O_z$, its distance from $I_z$ is at least $cB_\nu$, and its distance from $I_w$ is at least $c\rho$; the analogous bounds hold for exits of $O_w$. Equation~\eqref{eq:2.6} therefore bounds the two Green functions by $C+g\log(\rho_j/B_\nu)$ and $C+g\log(\rho_j/\rho)$, respectively. The capacities in~\eqref{eq:8.15} are bounded. Consequently, the uniform future-root mean satisfies
\begin{equation}\label{eq:8.17}
\mathcal R:=\sup_{y\in\mathcal S}\mathbb E_yK_y
\le C\left(1+\log\frac{\rho_j}{B_{\nu}}+\log\frac{\rho_j}{\rho}\right)
\le C(J+1).
\end{equation}

For $p\ge0$, put $f_p(y)=\mathbb E_y\binom{K_y}{p}$, with $f_0=1$. Let $T_i$ be the successive completion times of the future short roots. Grouping increasing-time selections by their first selected root and applying the strong Markov property at $T_i$ yields
\begin{equation}\label{eq:completion-recursion}
f_p(y)=\mathbb E_y\sum_{i\ge1}
\mathbf1_{\{T_i<\tau_{\Omega_j}\}}f_{p-1}(S_{T_i}),
\qquad p\ge1.
\end{equation}
Since every $S_{T_i}$ belongs to $\mathcal S$, induction gives $\sup_{y\in\mathcal S}f_p(y)\le\mathcal R^p$. Applying the same identity under the initial common-mark law, whose mean number of roots is $\mu$ from~\eqref{eq:8.16}, gives, for $K=K_z+K_w$,
\begin{equation}\label{eq:8.18}
\mathbf E\binom Kp\le\mu\mathcal R^{p-1},\qquad p\ge1.
\end{equation}
The sums are nonnegative, so the induction is justified by monotone convergence and also proves finiteness at each order.

For $W\ge0$ with $\mathcal RW<1$, the binomial expansion now gives
\begin{equation}\label{eq:8.19}
\mathbf E(1+W)^K
=1+\sum_{p\ge1}W^p\mathbf E\binom Kp
\le1+\frac{\mu W}{1-\mathcal RW}.
\end{equation}
Take $W=e^\xi-1$ with $\xi(J+1)$ sufficiently small. Using $e^\xi-1=\xi+O(\xi^2)$ and $\log(1+x)\le x$ gives $\log\mathbf E e^{\xi K}\le\mu\xi+C(J+1)\xi^2$. Independence of different common marks therefore implies, for $k$ marks,
\begin{equation}\label{eq:8.20}
\log\mathbf E_k e^{\xi(K_z+K_w)}
\le2k\xi+Ck(J+1)\xi^2+o(1).
\end{equation}
Here the error includes $k\xi(\mu-2)$. In the application $k=O(j^2)$ and $\xi=O(j^{-1})$, so~\eqref{eq:8.16} makes it uniformly $o(1)$ over the moving range.

\subsection{The moving early-distance sum}\label{sec:earlysum}
We now return to the exponential moment in~\eqref{eq:8.11}. Set $b=4+\theta j-j^\gamma$, the barrier height at the common endpoint. In this range $w_j=j^\gamma$ and $b\asymp j$. Concavity gives, for $K\ge0$,
\begin{equation}\label{eq:8.21}
2\sqrt{2K}\le b+2K/b.
\end{equation}
Apply this to each root and take $\xi=2/b$. Increasing the fixed $a_\circ$ makes $\xi(J+1)$ lie in the convergence range of~\eqref{eq:8.19}. Conditional on $u=\sqrt{2N_j}\le b$,~\eqref{eq:8.20} yields
\begin{equation}\label{eq:8.22}
\mathbf E_k e^{2\sqrt{2K_z}+2\sqrt{2K_w}}
\le\exp\{2b+2u^2/b+C(J+1)+o(1)\}.
\end{equation}

The geometric Galton--Watson process started from one individual has the exact marginal law
\begin{equation}\label{eq:8.23}
\mathbf P_1(Z_j=k)=\frac{j^{k-1}}{(j+1)^{k+1}},\qquad k\ge1.
\end{equation}
This follows by expanding its iterated generating function
$1-(1-z)/(1+j(1-z))$. In particular, the probability is bounded by
$e^{-u^2/[2(j+1)]}$. Since $2/b-1/[2(j+1)]>0$, the sum over
$1\le k\le b^2/2$ is bounded by the largest endpoint $u=b$, multiplied by $Cj^2$. Combining with~\eqref{eq:8.11}--\eqref{eq:8.12}, we obtain
\begin{equation}\label{eq:8.24}
\mathbb P_0(\mathcal C_z\cap\mathcal C_w)
\le C(\nu+1)^3j^2\frac{1}{R^4}
\exp\left[-4\nu-\frac{b^2}{2(j+1)}+4b+C(J+1)\right].
\end{equation}

Because $\theta\le2$, we have $b\le2j-j^\gamma+4$, and the function
$-b^2/[2(j+1)]+4b$ is increasing over this range. For $w=j^\gamma$,
\begin{equation}\label{eq:8.25}
-\frac{(2j-w+4)^2}{2(j+1)}+4(2j-w+4)
=6j-2w+10-\frac{(w-2)^2}{2(j+1)}.
\end{equation}
The number of site pairs with distance index $a$ is at most
\begin{equation}\label{eq:8.26}
CR^4e^{-2a}.
\end{equation}
By~\eqref{eq:8.8}, $-2a-4\nu+6j=-10J-20$. The remaining factor $e^{C(J+1)}$ is a fixed power of $a+1$. Thus, for some fixed finite $P$,
\begin{equation}\label{eq:8.27}
\sum_{a(z,w)=a}\mathbb P_0(\mathcal C_z\cap\mathcal C_w)
\le C(a+1)^P e^{-c a^{2/5}}.
\end{equation}
The sum over all moving early distances is uniformly bounded.

\subsection{Fixed early distances and macroscopic pairs}\label{sec:far}
For $a(z,w)\le a_\circ$, choose a fixed sufficiently large $J$ and then a fixed $\nu$ with $\nu-J\ge a_\circ+6$. The long balls satisfy the geometry in Section~\ref{sec:exterior}. The bound $K_\nu^{\max}$ and the comparison error are now fixed, and Lemma~\ref{lem:continuation} gives
$p_\nu(k)\le C_{a_\circ}R^{-2}$ for every allowed $k$. Apply~\eqref{eq:8.6} with $G$ the whole space to get
\begin{equation}\label{eq:8.28}
\mathbb P_0(\mathcal C_z\cap\mathcal C_w)
\le C_{a_\circ}R^{-4}.
\end{equation}
There are at most $CR^4$ such pairs. Their total contribution is bounded. This includes pairs farther apart than $\rho_0$, for which the index $a$ can be negative.


\section{Proof of Proposition~\ref{prop:seed}}\label{sec:second}
Recall that $W_R$, defined in~\eqref{eq:5.15}, counts successful sites and that $W_R>0$ implies $\sqrt{M_R}\ge b_R$. Proposition~\ref{prop:onepoint} bounds its first moment above and away from zero, and Section~\ref{sec:tworoot} bounds the contribution of early separations to its second moment. We estimate the remaining pairs, verify that the distance ranges cover every pair, and deduce the required positive probability.

\subsection{Pairs at bounded distance}\label{sec:neighbors}
For each site, there are $O_A(1)$ sites at distance at most $A$, including itself. Therefore, by~\eqref{eq:7.28},
\begin{equation}\label{eq:9.1}
\sum_{z,w:\,|z-w|\le A}
\mathbb P_0(\mathcal C_z\cap\mathcal C_w)
\le C_A\sum_z\mathbb P_0(\mathcal C_z)\le C_A.
\end{equation}
Since $A$ is fixed, the bound is uniform in $R$.

\subsection{Pairs within the terminal observation ball}\label{sec:latepairs}
Suppose $A<\rho=|z-w|\le a_{\max}/(8B)$, and choose the first observation radius $a_i\ge4\rho$, so that $4\rho\le a_i\le4B\rho$. Retain the coarse record and terminal band for $z$, the high-site conditions at both sites, and the auxiliary average path for $z$ up to time $v=v_i$. By~\eqref{eq:2.21}, all occupations used here come from the same collection of terminal marks. Normalize both $F_z$ and $F_w$ using this collection's fixed $k,u$ and intensity $\alpha$; the two high-site conditions then have the same threshold $A_t(k)$. Put
\[
 h=d-v=\log\rho+O(1),\qquad f=h^\gamma.
\]
We first suppose $h\le d/2$. The continuous average restriction implies
\begin{equation}\label{eq:pair-prefix}
 \mathcal B_z(s)\le q\quad(0\le s\le v),\qquad
 \mathcal B_z(v)\le q-f.
\end{equation}
The path restriction in~\eqref{eq:pair-prefix} supplies the factor $d^{-1}$ needed in the pair estimate, through the killed Brownian density below.

For the Gaussian vector consisting of $F_z,F_w,S_i,\ldots,S_I$, compensate each average variance by adding $\varepsilon_l$. Its centered covariance then has the exact representation
\begin{equation}\label{eq:pair-brownian}
 S_l-\mathbb ES_l=B(v_l),\quad
 F_z-\mathbb EF_z=B(v)+\xi,\quad
 F_w-\mathbb EF_w=B(v)+\zeta,
\end{equation}
where $B$ is Brownian motion and $(\xi,\zeta)$ is independent of $B$, with covariance $R_h=hI_2+O(1)$ entrywise. Indeed, the equilibrium identities give $\Cov(F_z,S_l)=\Cov(F_w,S_l)=v_l$ and the Brownian off-diagonal average covariances; the Green estimate gives $\Var F_z=\Var F_w=d+O(1)$ and $\Cov(F_z,F_w)=v+O(1)$. Taking the fixed $A$ sufficiently large gives $R_h\ge chI_2$ and hence a joint residual density bounded by $C/h$.

To compare with this Brownian representation uniformly in the number of averages, let $K$ be the covariance in~\eqref{eq:pair-brownian} and $E$ the diagonal matrix of average corrections. The quadratic form of $K^{-1}$ is the Brownian increment energy plus
\[
 (x_z-x_i,x_w-x_i)R_h^{-1}(x_z-x_i,x_w-x_i)^T.
\]
All Brownian time gaps are bounded below, so $\|K^{-1}\|\le C$. Also $\operatorname{tr}E\le C/A$. Thus $\|E^{1/2}K^{-1}E^{1/2}\|\le1/2$ for large $A$, and the exact density ratio yields
\begin{equation}\label{eq:pair-density}
 f_{K-E}(x)\le\det(I-E^{1/2}K^{-1}E^{1/2})^{-1/2}f_K(x)
 \le C f_K(x).
\end{equation}
It also gives a uniform lower eigenvalue bound for $K-E$.

Under the reference law with $k$ marks, use the exponential tilt $X=2F_z+2F_w-2S_i$. The mark moment expansion gives
\begin{equation}\label{eq:9.2}
 \Lambda_X:=\log\mathbf E_ke^X=4d-2v+O(1),
\end{equation}
where the error includes the third-order term $O(d^3/t)$. The leading tilted means are $2d+O(1)$ at both sites, $2v_l$ at each strictly outer average, and $2v+2\varepsilon_i$ at the last average. The tilt coefficients have absolute sum at most six. Together with the lower spectral bound above, this verifies the hypotheses of Lemma~\ref{lem:bounded-tilt} and gives the combined convex-set approximation error
\begin{equation}\label{eq:pair-clt}
 \delta_t\le Cd^5/t=o(d^{-2}).
\end{equation}
The signed form of that lemma applies to the negative coefficient of $S_i$.

We apply this absolute error bound after tilting and sum it with the inverse-tilt weights. Split the site excesses and the last average deficit into unit windows indexed by $n,m,l\ge0$:
\begin{equation}\label{eq:pair-windows}
\begin{gathered}
 F_z\in[A_t(k)+n,A_t(k)+n+1),\quad
 F_w\in[A_t(k)+m,A_t(k)+m+1),\\
 S_i\in[2v+q-f-l-1,2v+q-f-l).
\end{gathered}
\end{equation}
Include the possible top boundary in the first deficit window. With the auxiliary bridges fixed, these windows together with~\eqref{eq:pair-prefix} define a convex set of observation vectors; the half-open conventions follow by taking limits of convex sets. In the compensated Gaussian law, the mean correction is zero before the last retained average interval and interpolates linearly from zero to $2\varepsilon_i$ on that interval. Its nonnegativity tightens the flat Brownian constraint. For deficit $s\ge0$, with $p_v(x)=(2\pi v)^{-1/2}e^{-x^2/(2v)}$, the killed Brownian density satisfies
\[
 p_v(q-s)-p_v(q+s)\le C\frac{qs}{v^{3/2}}.
\]
In the last-average window the deficit belongs to
$[f+l+2\varepsilon_i,f+l+2\varepsilon_i+1]$. Since $v\ge d/2$ and $q\asymp\sqrt d$, its endpoint probability is at most $C(f+l+1)/d$.
Multiplication by the residual density $C/h$ bounds the Gaussian probability of each window by $C(f+l+1)/(dh)$.

On that window,~\eqref{eq:7.19} and~\eqref{eq:9.2} bound the inverse tilt by
\[
 e^{\Lambda_X-X}\le
 C e^{-2d-2h-2q-2f}e^{-2(n+m+l)}.
\]
Apply~\eqref{eq:pair-density}, add the approximation error~\eqref{eq:pair-clt} in each convex window, and sum against the three geometric weights. This gives the discrete mark estimate
\begin{equation}\label{eq:9.3}
 \mathbf P_k(F_z,F_w\ge A_t(k),\ \eqref{eq:pair-prefix})
 \le C\frac{1+f}{dh}e^{-2d-2h-2q-2f}.
\end{equation}
The summed absolute error is at most $C\delta_t$ times the displayed exponential prefactor. Since $\delta_t=o(d^{-2})$, it is absorbed uniformly for $h\le d/2$. Thus the convex-set approximation, combined with the inverse-tilt weights, has sufficient precision for~\eqref{eq:9.3}.

Multiplying~\eqref{eq:9.3} by the coarse-window estimate~\eqref{eq:7.24}, applying the original-record comparison~\eqref{eq:5.17}, and summing the $q$ band gives
\begin{equation}\label{eq:9.4}
 \mathbb P_0(\mathcal C_z\cap\mathcal C_w)
 \le CR^{-2}\rho^{-2}\frac{1+h^\gamma}{h}e^{-c h^\gamma}.
\end{equation}
There are at most $CR^2a_i^2$ ordered site pairs assigned to $a_i$. Since the values $h_i=d-v_i$ are separated by fixed positive constants, their shell contributions are bounded by the convergent series
$C\sum_i(1+h_i^\gamma)h_i^{-1}e^{-c h_i^\gamma}$.
For $h>d/2$, Chernoff's inequality alone, using the same signed statistic and the endpoint restriction, gives~\eqref{eq:9.3} without its prefactor $(1+f)/(dh)$. After the coarse sum and spatial counting, each shell costs at most $Cd e^{-2f}$; there are $O(d)$ shells and $f\ge(d/2)^\gamma$. We have therefore proved
\begin{equation}\label{eq:9.5}
 \sum_{A<|z-w|\le a_{\max}/(8B)}
 \mathbb P_0(\mathcal C_z\cap\mathcal C_w)
 \le C+Cd^2e^{-c d^\gamma}\le C.
\end{equation}

\subsection{Pairs with a later coarse common endpoint}\label{sec:coarsepairs}
For the distance index in~\eqref{eq:8.7}, let
\begin{equation}\label{eq:9.6}
J_t=\lceil10\log t\rceil+5,\qquad j_t=a(z,w)-J_t.
\end{equation}
For $\mathcal L_t\le j=j_t\le N_R$, both sites can be estimated within a single collection of common marks. Use the coarse prefix centered at $z$ and observe a ball whose radius is a fixed multiple of $\rho$ and which contains both sites. Since $\rho/\rho_{j+1}=O(e^{-J_t})$, the logarithmic comparison error is
\begin{equation}\label{eq:9.7}
Ct^2\{H_0/\rho_{j+1}+(J_t+1)e^{-J_t}\}=o(1).
\end{equation}
We retain both site occupations in this common mark law.

In the common killing disk $B(z,\rho_j+H_0)$, the two-site Green function satisfies
\[
G_{zz}/g=t-j+O(1),\qquad G_{ww}/g=t-j+O(1),\qquad
G_{zw}/g=J_t+O(1).
\]
Put
\begin{equation}\label{eq:9.8}
D_j=g^{-1}\max_{v\in\{z,w\}}
\sum_{v'\in\{z,w\}}(G-I)_{vv'}
=t-j+J_t+O(1),
\end{equation}
The mark intensity is $\alpha_j$ from~\eqref{eq:8.13}. For $u=\sqrt{2N_j}$, the nonnegative resolvent series in~\eqref{eq:2.16} and the row-sum bound give, whenever $W>0$ and $gD_jW<1$,
\begin{equation}\label{eq:9.9}
\mathbf P_{u^2/2}^{(j)}(H_z,H_w\ge k_R)
\le\exp\left\{-gy^2\log(1+W)
+\frac{\alpha_j u^2W}{1-gD_jW}\right\}.
\end{equation}
The superscript specifies the common mark law at level $j$, and the subscript gives its fixed number of marks. In particular, $D_j$ uses the row sums of the discrete kernel $G-I$ appearing in~\eqref{eq:2.16}.

Let $b=4+\theta j-w_j$ and $\bar b=\sqrt{\alpha_j/g}\,b$. For every permitted endpoint $u\le b$, use the same parameter
\begin{equation}\label{eq:9.10}
W_b=\frac{1-\bar b/y}{gD_j}=O(1/t).
\end{equation}
Since $T-j\ge d\gg J_t$, this parameter satisfies $0<gD_jW_b<1$. Using $\log(1+W)\ge W-W^2/2$ costs $gy^2W_b^2=O(1)$. After multiplication by the common geometric GW marginal~\eqref{eq:8.23}, the coefficient of $u^2$ is
\[
\frac{\alpha_jW_b}{1-gD_jW_b}-\frac1{2(j+1)}>0.
\]
Indeed, its first term is $(1+o(1))/j$, uniformly in the range, because
$w_j/j$, $w_j/(T-j)$ and $J_t/(T-j)$ all tend to zero. The largest allowed endpoint is therefore $u=b$. Summing at most $Ct^2$ endpoints and using $\bar b-b=o(1)$ gives
\begin{equation}\label{eq:9.11}
\mathbb P_0(\text{common barrier, both sites high})
\le Ct^2\exp\left[-\frac{b^2}{2(j+1)}-\frac{(y-b)^2}{D_j}\right].
\end{equation}

To compare this exponential bound with the number of site pairs, write $b=\theta j+v$, where $v=4-w_j$. With denominators $j$ and $T-j$, the identity is
\begin{equation}\label{eq:9.12}
\frac{b^2}{2j}+\frac{(y-b)^2}{T-j}
=\theta^2(T-j/2)-\theta v
+v^2\left(\frac1{2j}+\frac1{T-j}\right).
\end{equation}
This is at least $4t-2j-4\log t+2w_j-C$. Replacing $j$ by $j+1$ loses at most a constant. Replacing $T-j$ by
$D_j=T-j+J_t+O(1)$ loses at most $4J_t+C$, since
\[
\frac{y-b}{T-j}=\theta+O\bigl(w_j/(T-j)\bigr),
\qquad J_t w_j/(T-j)\le C J_t d^{\gamma-1}=o(1).
\]
Also $(\theta-2)w_j=o(1)$ uniformly. Hence
\begin{equation}\label{eq:9.13}
\frac{b^2}{2(j+1)}+\frac{(y-b)^2}{D_j}
\ge4t-2j+2w_j-4\log t-4J_t-C.
\end{equation}

There are at most $CR^4e^{-2(j+J_t)}$ ordered site pairs at this common endpoint index. Equations~\eqref{eq:9.11}--\eqref{eq:9.13} yield
\begin{equation}\label{eq:9.14}
\sum_{j_t(z,w)=j}\mathbb P_0(\mathcal C_z\cap\mathcal C_w)
\le Ct^{26}e^{-2w_j}.
\end{equation}
Since $j\ge\mathcal L_t$ and $T-j\ge d$, one has
$w_j\ge1000\log t+o(1)$. The sum over at most $t$ levels tends to zero.

\subsection{Partition of all site pairs}\label{sec:partition}
We now combine the estimates, assigning each ordered pair to the first applicable range. The diagonal and pairs with $\rho\le A$ satisfy~\eqref{eq:9.1}. Pairs with $A<\rho\le a_{\max}/(8B)$ satisfy~\eqref{eq:9.5}. Among the remaining pairs, those with $j_t\ge\mathcal L_t$ satisfy~\eqref{eq:9.14}. The pairs left over satisfy~\eqref{eq:8.28} if $a\le a_\circ$, and~\eqref{eq:8.27} if $a>a_\circ$.

It remains to check the generation bounds in these ranges. If $\rho>a_{\max}/(8B)$, then $a_{\max}>m/(Bt^6)$ and $\rho_0=me^{N_R+1}$ give
\begin{equation}\label{eq:9.15}
j_t\le N_R+1+6\log t+\log(8B^2)-J_t
\le N_R-4\log t+C_B<N_R.
\end{equation}
Thus the generation in the coarse-pair estimate is at most $N_R$. If a remaining pair has $j_t<\mathcal L_t$, then
\begin{equation}\label{eq:9.16}
a<\mathcal L_t+J_t,
\end{equation}
as required for the early-pair estimates in Section~\ref{sec:tworoot}. The choice $J_t$ applies to the later coarse range, while $J$ from~\eqref{eq:8.8} applies to the early range and varies with the distance index.

\subsection{Proof of Proposition~\ref{prop:seed}}\label{sec:seedcompletion}
Proposition~\ref{prop:onepoint} and the five distance contributions give, for all sufficiently large real $R$,
\begin{equation}\label{eq:9.17}
0<c_1\le\mathbb EW_R\le C_1,\qquad
\mathbb EW_R^2\le C_2.
\end{equation}
By Cauchy--Schwarz and the pathwise implication~\eqref{eq:5.16},
\begin{equation}\label{eq:9.18}
\mathbb P_0(\sqrt{M_R}\ge b_R)
\ge\mathbb P_0(W_R>0)
\ge\frac{(\mathbb EW_R)^2}{\mathbb EW_R^2}
\ge c_1^2/C_2=:p_*>0.
\end{equation}
This proves Proposition~\ref{prop:seed}. Section~\ref{sec:mainproofs} then gives both main theorems. The argument applies to all sufficiently large real $R$, since rounding affects only the bounded difference $T-t$ and the one-step overshoot estimates.
\qed

\appendix
\section{Formalization and correspondence}\label{sec:formalization}
The accompanying Lean development records the statements of this manuscript in a label-to-declaration correspondence. Its main declaration is
\begin{center}
\texttt{source\_main\_claims\_internal\_reflection}.
\end{center}
It proves the disk and deterministic-time tightness and tail claims, including the inclusive-time tightness variant. This is a conditional formalization with six explicit literature input families: the potential-kernel asymptotic and positive-harmonic relative estimates from~\cite{LL}, the upper and lower critical geometric Galton--Watson barrier estimates from~\cite{GW}, the convex-set normal approximation from~\cite{BE}, and the normalized dyadic-square DGFF maximum estimate from~\cite{BDZ}. These inputs remain parameters of the Lean theorem; they are not counted as internally proved results.

The correspondence concerns the mathematical objects, assumptions, parameter ranges and conclusions. Some auxiliary arguments use different routes. For~\eqref{eq:2.5}, this manuscript cites the explicit potential-kernel constant, whereas Lean identifies it from diagonal potential values using an asymptotic input with an initially unspecified constant. The full-radius comparison in~\eqref{eq:2.24} uses the uniform entrance comparison and the radius split given in Sections~\ref{sec:marks} and~\ref{sec:recordcomparison}; these steps make the lower exponential bound valid even when the relative error is not small. The return estimate in Section~\ref{sec:roots} uses the buffer to absorb the boundary error uniformly in the root level. The finite-time killed-kernel derivation of~\eqref{eq:7.17} follows the route used in Lean.

The machine-readable correspondence is maintained in \texttt{obligations.json}. It distinguishes internally proved statements, statements conditional on the listed inputs, and explicitly external statements. Formal acceptance of a declaration verifies its derivation from its stated assumptions; the correspondence with the manuscript requires a separate check of definitions and scope. The source checksum and correspondence records must therefore be renewed together after a manuscript revision.

\clearpage
\phantomsection
\addcontentsline{toc}{section}{References}
\begin{thebibliography}{99}
\raggedright

\bibitem{BDZ}
M.~Bramson, J.~Ding and O.~Zeitouni.
Convergence in law of the maximum of the two-dimensional discrete Gaussian free field.
\emph{Communications on Pure and Applied Mathematics} \textbf{69}(1) (2016), 62--123.
The estimates cited here are numbered as in arXiv:1301.6669v5.

\bibitem{DPRZ}
A.~Dembo, Y.~Peres, J.~Rosen and O.~Zeitouni.
Thick points for planar Brownian motion and the Erd\H{o}s--Taylor conjecture on random walk.
\emph{Acta Mathematica} \textbf{186}(2) (2001), 239--270.

\bibitem{LL}
G.~F.~Lawler and V.~Limic.
\emph{Random Walk: A Modern Introduction}.
Cambridge Studies in Advanced Mathematics, vol.~123.
Cambridge University Press, Cambridge, 2010.

\bibitem{BE}
M.~Rai\v{c}.
A multivariate Berry--Esseen theorem with explicit constants.
\emph{Bernoulli} \textbf{25}(4A) (2019), 2824--2853.
See also arXiv:1802.06475v4.

\bibitem{ET}
J.~Rosen.
A random walk proof of the Erd\H{o}s--Taylor conjecture.
\emph{Periodica Mathematica Hungarica} \textbf{50} (2005), 223--245.
The theorem numbering used here is that of arXiv:math/0503108v1.

\bibitem{GW}
J.~Rosen.
Tightness for thick points in two dimensions.
\emph{Electronic Journal of Probability} \textbf{28} (2023), paper no.~18, 1--45.
The theorem and section numbering used here is that of arXiv:2203.14394v1.

\end{thebibliography}

\end{document}
